%% Delete this \nocite command invocation to make the references section only list out the bibitems that are actually cited.
\nocite{*}


% ============================================================
% Section 1: Algebraic Preliminaries
% ============================================================
\section{Algebraic Preliminaries}

\subsection{Groups}

\begin{lemma} \label{lemma:cancellation_law_for_group_operations}
    Let $(G,\cdot)$ be a \CrefAndHyperrefIfExist{definition:group}{group}. For any $a,b,c \in G$,
    \begin{enumerate}
        \item $ab = ac$ if and only if $b = c$.
        \item $ba = bc$ if and only if $a = c$.
    \end{enumerate}
\end{lemma}
\begin{proof}
    \begin{enumerate}
        \item The ``if'' direction is trivial. To prove the ``only if'' direction, suppose thta $a,b,c \in G$ satisfy $ab = ac$. Applying $a^{-1}$ to both sides, we have
        $$a^{-1} \cdot (ab) = a^{-1} \cdot (ac).$$
        By the associativity of $\cdot$, this equality becomes
        $$(a^{-1} a) b = (a^{-1} \cdot a)c,$$
        which in turn becomes
        $$e b = ec.$$
        Since $e$ is the identity element, we have $b = c$ as desired.
        \item One obtains this through a symmetric argument.
    \end{enumerate}
\end{proof}

\begin{corollary} \label{corollary:cancellation_law_for_identity_element_of_group_operation}
    Let $(G,\cdot)$ be a \CrefAndHyperrefIfExist{definition:group}{group}. For any $a,b \in G$,
    \begin{enumerate}
        \item if $ab = b$, then $a = e$.
        \item if $ab = a$, then $b = e$.
    \end{enumerate}
\end{corollary}
\begin{proof}
    \begin{enumerate}
        \item Let $a,b \in R$ satisfy $ab = b$. In particular, $ab = eb$. Apply \Cref{lemma:cancellation_law_for_group_operations}.

        \item One obtains this through a symmetric argument.
    \end{enumerate}
\end{proof}

\begin{corollary} \label{corollary:elements_of_group_that_multiply_to_identity_are_inverses}
    Let $(G,\cdot)$ be a \CrefAndHyperrefIfExist{definition:group}{group}. For any $a,b \in G$, if $ab = e$, then $b = a^{-1}$, and $a = b^{-1}$.
\end{corollary}
\begin{proof}
    Let $a,b \in G$ satisfy $ab = e$. Since $e = a \cdot a^{-1}$, we have
    $$a \cdot b = a \cdot a^{-1}.$$
    Apply \Cref{lemma:cancellation_law_for_group_operations} to obtain that $b = a^{-1}$. Similarly, $e = b^{-1} \cdot b$, so
    $$a \cdot b  = b^{-1} \cdot b$$
    and hence $a = b^{-1}$ by \Cref{lemma:cancellation_law_for_group_operations}.
\end{proof}

\begin{corollary} \label{corollary:inverse_of_inverse_in_group_is_element}
    Let $(G,\cdot)$ be a \CrefAndHyperrefIfExist{definition:group}{group}. For any $a \in G$, $\left( a^{-1} \right)^{-1} = a$.
\end{corollary}
\begin{proof}
    Note that $\left( a^{-1} \right) \cdot a = e$. Applying \Cref{corollary:elements_of_group_that_multiply_to_identity_are_inverses}, we have that $\left( a^{-1} \right)^{-1} = a$. 
\end{proof}

\subsection{Rings}

\input{../_definitions/definition_ring.tex}
\begin{lemma} \label{lemma:multiplication_by_zero_in_a_ring_is_zero}
    Let $R$ be a \CrefAndHyperrefIfExist{definition:ring}{ring}. For any $r \in R$, we have $0 \cdot r = 0 = r \cdot 0$.
\end{lemma}
\begin{proof}
    For any $r \in R$, we have 
    $$0 \cdot r = (0+0) \cdot r = 0 \cdot r + 0 \cdot r.$$
    By \Cref{corollary:cancellation_law_for_identity_element_of_group_operation}, we have that $0 = 0 \cdot r$. Similarly, $r \cdot 0 = 0$ as well.
\end{proof}
\begin{lemma} \label{lemma:multiplication_by_minus_one_in_a_ring_is_additive_inverse}
    Let $R$ be a \CrefAndHyperrefIfExist{definition:ring}{ring}. For any $r \in R$, we have $(-1) \cdot r = -r = r \cdot (-1)$.
\end{lemma}
\begin{proof}
    For any $r \in R$, distributivity of $\cdot$ over $+$ yields
    $$(1 + (-1)) \cdot r = 1 \cdot r + (-1) \cdot r.$$
    On the one hand,
    $$(1 + (-1)) \cdot r = 0 \cdot r = 0$$
    by \Cref{lemma:multiplication_by_zero_in_a_ring_is_zero}. 
    On the other hand,
    $$1 \cdot r + (-1) \cdot r = r + (-1) \cdot r,$$
    so 
    $$ 0 = r+ (-1) \cdot r.$$
    By \Cref{corollary:elements_of_group_that_multiply_to_identity_are_inverses}, 
    $$(-1) \cdot r = -r.$$
    Similarly, recognizing that 
    $$((-1) + 1) \cdot r = (-1) \cdot r + 1 \cdot r,$$
    one can formulate a similar argument to show that $r \cdot (-1) = -r$. 
\end{proof}
\begin{corollary} \label{corollary:minus_one_times_minus_one_is_one_in_a_ring}
    Let $R$ be a \CrefAndHyperrefIfExist{definition:ring}{ring}. We have 
    $$(-1) \cdot (-1) = 1.$$
\end{corollary}
\begin{proof}
    By \Cref{lemma:multiplication_by_minus_one_in_a_ring_is_additive_inverse}, we have
    $$(-1) \cdot (-1) = -(-1).$$
    By \Cref{corollary:inverse_of_inverse_in_group_is_element}, $-(-1)$ equals $1$.
\end{proof}


% ============================================================
% Section 2: Natural Numbers
% ============================================================
\section{Natural Numbers}

\subsection{Definition and Basic Properties}

\begin{definition} \label{definition:natural_numbers}
Let $\bbN$ be a set. The \hldef{set of natural numbers}, $\bbN$, is equipped with:
\begin{enumerate}
    \item a distinguished element \hl{$1 \in \bbN$}, called the \hldef{first natural number}; and
    \item a function \hl{$s: \bbN \to \bbN$}, called the \hldef{successor function},
\end{enumerate}
such that the following axioms, called the \hldef{Peano axioms}, hold:
\begin{enumerate}[resume]
    \item $s$ is \CrefAndHyperrefIfExist{definition:injective_surjective_bijective_map_of_sets}{injective}.
    \item $1$ is not in the image of $s$.
    \item (\hldef{Principle of mathematical induction}) If $S \subseteq \bbN$ is a subset such that
    \begin{enumerate}
        \item $1 \in S$, and
        \item for every $n \in S$, we have $s(n) \in S$,
    \end{enumerate}
    then $S = \bbN$.
\end{enumerate}

The successor of $n$ is written \hl{$s(n)$} or \hl{$n+1$}\TextIfExists{proposition:natural_numbers_addition_existence}{ (see also \Cref{proposition:natural_numbers_addition_existence})}.
\end{definition}

\begin{remark} \label{remark:natural_numbers_zero_inclusive}
Some texts include $0$ in the natural numbers. In this project, when $0$ is included we explicitly write $\bbN_0 = \bbN \cup \{0\}$.
\end{remark}

\begin{remark} \label{remark:natural_numbers_zero_inclusive}
Some texts include $0$ in the natural numbers. In this project, when $0$ is included we explicitly write $\bbN_0 = \bbN \cup \{0\}$.
\end{remark}
\begin{proposition} \label{proposition:natural_numbers_addition_existence}
Let $(\bbN, 1, s)$ be the set of \CrefAndHyperrefIfExist{definition:natural_numbers}{natural numbers}. For each $m \in \bbN$, let $f_m: \bbN \to \bbN$ be the unique recursive \CrefAndHyperrefIfExist{definition:function_of_sets}{function} on the initial value $s(m)$ and the recursive relation $s: \bbN \to \bbN$ whose existence and uniqueness is guaranteed by \Cref{proposition:recursive_function_existence_uniqueness}. In other words,
\begin{align*}
f_m(1) &= s(m) \\
f_m(s(n)) &= s(f_m(n)) \text{ for } n \in \bbN.
\end{align*}

Define
\hl{$n + m := f_m(n)$.}
Then $+: \bbN \times \bbN \to \bbN$ is a \CrefAndHyperrefIfExist{definition:n_ary_operation_binary_operation_ternary_operation_on_a_set}{binary operation} satisfying for all $n, m \in \bbN$:
\begin{enumerate}
    \item $n + 1 = s(n)$,
    \item $n + s(m) = s(n + m)$.
\end{enumerate}
This binary operation is called \hldef{addition in $\bbN$}.
\end{proposition}

\begin{proof}
To show that $n+1 = s(n)$ for all $n \in \bbN$ is to show that $f_1(n) = s(n)$ for all $n \in \bbN$. We do so by induction. For $n = 1$, $f_1(1) = s(1)$ by the defining property of $f_1$. Now suppose inductively that $f_1(k) = s(k)$ for some $k \in \bbN$. It follows that $s(f_1(k)) = s(s(k))$. The left hand side equals $f_1(s(k))$, so $f_1(s(k)) = s(s(k))$. Therefore, $n+1 = s(n)$ for all $n \in \bbN$ by induction.

To show that $n+s(m) = s(n+m)$ for all $n,m \in \bbN$, fix $m \in \bbN$. 

Consider the function $h: \bbN \to \bbN$ defined by $h(n) = n + s(m) = f_{s(m)}(n)$. By definition of $f_{s(m)}$, we have $h(1) = f_{s(m)}(1) = s(s(m))$ and $h(s(k)) = s(h(k))$ for all $k$. On the other hand, consider the function $\bbN \to \bbN, \quad k \mapsto s(k + m) = s(f_m(k))$. This function 
\begin{enumerate}
    \item sends $k = 1$ to $s(f_m(1)) = s(s(m))$, and 
    \item sends $k = s(j)$ to $s(f_{m}(s(j))) = s(s(f_m(j)))$.
\end{enumerate}
 Thus both functions $n \mapsto n + s(m)$ and $n \mapsto s(n + m)$ are recursive functions on the same pair $(s(s(m)), s)$. Recurisve functions are unique by \CrefAndHyperrefIfExist{proposition:recursive_function_existence_uniqueness}{the recursive function theorem}, so they coincide, i.e. $n + s(m) = s(n + m)$ for all $n$.
\end{proof}

\begin{proposition} \label{proposition:natural_numbers_multiplication_existence}
Let $(\bbN, 1, s)$ be the set of \CrefAndHyperrefIfExist{definition:natural_numbers}{natural numbers} and let $+: \bbN \times \bbN \to \bbN$ be \CrefAndHyperref{proposition:natural_numbers_addition_existence}{addition}. For each $m \in \bbN$, let $g_m: \bbN \to \bbN$ be the unique recursive \CrefAndHyperrefIfExist{definition:function_of_sets}{function} on the initial value $m$ and the recursive relation $(\cdot) + m : \bbN \to \bbN, x \mapsto x + m$, whose existence and uniqueness is guaranteed by \Cref{proposition:recursive_function_existence_uniqueness}. In other words,
\begin{align*}
	g_m(1) &= m \\
	g_m(s(n)) &= g_m(n) + m \text{ for } n \in \bbN.
\end{align*}

Define
\hl{$n \times m := g_m(n)$.}
Then $\times: \bbN \times \bbN \to \bbN$, also denoted \hl{$\cdot$}, is a \CrefAndHyperrefIfExist{definition:n_ary_operation_binary_operation_ternary_operation_on_a_set}{binary operation} satisfying for all $n, m \in \bbN$:
\begin{enumerate}
    \item $n \times 1 = n$,
    \item $n \times s(m) = n \times m + n$.
\end{enumerate}
\CrefIfExists{proposition:recursive_function_existence_uniqueness}
This binary operation is called \hldef{multiplication in $\bbN$}. 

Multiplication is universally notated to take precedent over addition, unless parentheses are written, e.g. $n \times m + k$ is universally understood as $(n \times m) + k$ and not as $n \times (m+k)$.
\end{proposition}

\begin{proof}
To show that $n \times 1 = n$ for all $n \in \bbN$ is to show that $g_1(n) = n$ for all $n \in \bbN$. We do so by induction. For $n = 1$, $g_1(1) = 1$ by the defining property of $g_1$. Now suppose inductively that $g_1(k) = k$ for some $k \in \bbN$. It follows that $g_1(s(k)) = g_1(k) + 1 = k + 1 = s(k)$, where the last equality holds by \Cref{proposition:natural_numbers_addition_existence}. Therefore, $n \times 1 = n$ for all $n \in \bbN$ by induction.

To show that $n \times s(m) = n \times m + n$ for all $n, m \in \bbN$, fix $m \in \bbN$.

Consider the function $h: \bbN \to \bbN$ defined by $h(n) = n \times s(m) = g_{s(m)}(n)$. By the definition of $g_{s(m)}$, we have $h(1) = g_{s(m)}(1) = s(m)$ and $h(s(k)) = h(k) + s(m)$ for all $k$. On the other hand, consider the function $\bbN \to \bbN, \quad k \mapsto k \times m + k = g_m(k) + k$. This function
\begin{enumerate}
    \item sends $k = 1$ to $g_m(1) + 1 = m + 1 = s(m)$, and
    \item sends $k = s(j)$ to $g_m(s(j)) + s(j) = (g_m(j) + m) + s(j)$. By \CrefAndHyperrefIfExist{proposition:natural_numbers_addition_associative_commutative}{associativity of addition}, this equals $g_m(j) + (m + s(j))$. By \Cref{proposition:natural_numbers_addition_existence}, $m + s(j) = s(m + j)$; by \CrefAndHyperrefIfExist{proposition:natural_numbers_addition_associative_commutative}{commutativity of addition}, $s(m + j) = s(j + m)$; by \Cref{proposition:natural_numbers_addition_existence} again, $s(j + m) = j + s(m)$. Hence the expression equals $g_m(j) + (j + s(m)) = (g_m(j) + j) + s(m)$.
\end{enumerate}
Thus both functions $n \mapsto n \times s(m)$ and $n \mapsto n \times m + n$ are recursive functions on the same pair $(s(m), (\cdot) + s(m))$. Recursive functions are unique by \CrefAndHyperrefIfExist{proposition:recursive_function_existence_uniqueness}{the recursive function theorem}, so they coincide, i.e. $n \times s(m) = n \times m + n$ for all $n$.
\end{proof}

\begin{proposition} \label{proposition:natural_numbers_addition_associative_commutative}
Let $\bbN$ be the set of \CrefAndHyperrefIfExist{definition:natural_numbers}{natural numbers}. For all $n, m, k \in \bbN$:
\begin{enumerate}
    \item $(n + m) + k = n + (m + k)$ (associativity),
    \item $n + m = m + n$ (commutativity),
    \item If $n + k = m + k$, then $n = m$ (cancellativity).
\end{enumerate}
\CrefIfExists{proposition:natural_numbers_addition_existence} 
\TextIfExists{definition:semigroup}{
In other words, $(\bbN,+)$ is a \CrefAndHyperrefIfExist{definition:semigroup}{commutative semigroup}.
}
\end{proposition}

\begin{proof}
    \begin{enumerate}
        \item Fix $n, m$ and induct on $k$. For $k = 1$, 
        $$(n + m) + 1 = s(n + m) = n + s(m) = n + (m + 1).$$ 
        If the identity holds for $k$, then 
        $$(n + m) + s(k) = s((n + m) + k) = s(n + (m + k)) = n + s(m + k) = n + (m + s(k)).$$

        \item We first prove that $1+n = n+1$ for all $n \in \bbN$ by induction on $n$. Clearly, $1+1 = 1+1$. Suppose inductively that $1+k = k+1$ for some $k \in \bbN$, i.e. $1+k = s(k)$. It follows that 
        \begin{align*}
        1+s(k) &= s(1+k) = s(k+1) = k + s(1) = k + (1+1) \\
        &= (k+1) +1 = (1+k) +1 = s(k) +1.
        \end{align*}
        Now we induct in $n$ for fixed $m$: the base case is true because we showed $1+m = m+1$ for any $m \in \bbN$. Suppose that $k \in \bbN$ satisfies $k+m = m+k$. By \CrefAndHyperrefIfExist{lemma:successor_left_addition}{the successor-left addition lemma}, $s(k)+m = s(k+m)$. Hence
        $$s(k) + m = s(k+m) = s(m+k) = m + s(k).$$

        \item Fix $n$ and $m$ and induct on $k$. For $k=1$, suppose $n+1 = m+1$. By property (i) of addition, $s(n) = s(m)$, so $n=m$ by injectivity of the successor function. Now suppose inductively that $k \in \bbN$ satisfies the implication $(n+k = m+k)  \Rightarrow (n = m)$. Suppose that $n+s(k) = m+s(k)$. It follows that $s(n+k) = s(m+k)$. By injectivity of the successor function, $n+k = m+k$, so by the inductive hypothesis, $n=m$.
    \end{enumerate}

\end{proof}

\begin{proposition} \label{proposition:natural_numbers_multiplication_associative_commutative_distributive}
Let $\bbN$ be the set of \CrefAndHyperrefIfExist{definition:natural_numbers}{natural numbers}. For all $n, m, k \in \bbN$:
\begin{enumerate}
    \item $(n \times m) \times k = n \times (m \times k)$ (associativity),
    \item $n \times m = m \times n$ (commutativity),
    \item $n \times (m + k) = n \times m + n \times k$ (distributivity).
\end{enumerate}
\CrefIfExists{proposition:natural_numbers_addition_existence}\CrefIfExists{proposition:natural_numbers_multiplication_existence}
Moreover, $1$ is the multiplicative identity: $n \times 1 = n$ for all $n \in \bbN$.
\end{proposition}

\begin{proof}

        We first show the distributivity property by fixing $n,m \in \bbN$ and inducting on $k$.  
        \begin{itemize}
            \item For the base case, i.e. $k = 1$, we have 
            $$n \times (m+1) = n \times s(m) = n \times m + n = n \times m + n \times 1$$
            by \Cref{proposition:natural_numbers_addition_existence} and \Cref{proposition:natural_numbers_multiplication_existence}.

            \item Suppose inductively that $l \in \bbN$ satisfies
            $$n \times (m+l) = n \times m + n \times l.$$
            We need to show that 
            $$n \times (m+s(l)) = n \times m + n \times s(l).$$
            The left hand side equals
            $$n \times s(m+l) = n \times (m+l) + n = (n \times m + n \times l) + n$$
            by \Cref{proposition:natural_numbers_addition_existence} and the inductive hypothesis. By the associtiativy property in \Cref{proposition:natural_numbers_addition_associative_commutative}, this equals 
            $$(n \times m) + (n \times l + n)$$
            which equals
            $$(n \times m) + (n \times s(l))$$
            as desired.
        \end{itemize}


        We show that $\times$ is associative by fixing $n, m$ and inducting on $k$.
        \begin{itemize}
            \item For the base case, i.e. $k = 1$, by \Cref{proposition:natural_numbers_multiplication_existence}, we have
            $$(n \times m) \times 1 = n \times m$$
            and
            $$n \times (m \times 1) = m \times m$$
            Therefore, 
            $$(n \times m) \times 1 = n \times (m \times 1).$$

            \item Suppose inductively that $l \in \bbN$ is an integer such that 
            $$(n \times m) \times l = n \times (m \times l);$$
            we aim to show that 
            $$(n \times m) \times s(l) = n \times (m \times s(l)).$$
            \CrefIfExists{definition:natural_numbers}
            By \Cref{proposition:natural_numbers_multiplication_existence}, 
            $$(n \times m) \times s(l) = (n \times m) \times l + (n \times m),$$
            which by the inductive hypothesis equals 
            $$n \times (m \times l) + n \times m.$$
            In turn, by distributivity, this equals 
            $$n \times (m \times l + m)$$
            and once again, by \Cref{proposition:natural_numbers_multiplication_existence}, this equals 
            $$n \times (m \times s(l)).$$
            We have thus shown the desired equality.
        \end{itemize}

        We show that $\times$ is commutative by inducting on $n$ for fixed $m$.

        \begin{itemize}
            \item For the base case, i.e. $n = 1$, $1 \times m = m$ by \Cref{lemma:natural_number_one_left_identity_for_multiplication}, and $m \times 1 = m$ by \CrefIfExists{proposition:natural_numbers_multiplication_existence}. Hence $1 \times m = m \times 1$.

            \item Suppose that $k \in \bbN$ satisfies $k \times m = m \times k$. By \Cref{lemma:natural_number_successor_left_multiplication_recursion}, the inductive hypothesis, and \CrefIfExists{proposition:natural_numbers_multiplication_existence}, we have
            $$s(k) \times m = k \times m + m = m \times k + m = m \times s(k).$$
        \end{itemize}

\end{proof}

\begin{definition} \label{definition:standard_order_on_natural_numbers}
Let $\bbN$ be the set of \CrefAndHyperrefIfExist{definition:natural_numbers}{natural numbers}, and let $+$ be \CrefAndHyperrefIfExist{proposition:natural_numbers_addition_existence}{addition in $\bbN$}. For $m, n \in \bbN$, we define the \hldef{standard order on $\bbN$} by declaring \hl{$m < n$} if there exists $k \in \bbN$ such that
$$m + k = n.$$
The \CrefAndHyperrefIfExist{definition:reflexive_closure_of_a_binary_relation_on_a_set}{reflexive closure} of this order is denoted \hl{$\leq$}. We write \hl{$n > m$} for $m < n$, and \hl{$n \geq m$} for $m \leq n$. Accordingly, to say that \hldef{$m$ is less than $n$} or that \hldef{$n$ is greater than $m$} is to say that $m < n$ or that $n > m$.
\end{definition}
\begin{proposition} \label{proposition:natural_numbers_order_compatibility}
Let $\bbN$ be the set of \CrefAndHyperrefIfExist{definition:natural_numbers}{natural numbers} with \CrefAndHyperrefIfExist{definition:standard_order_on_natural_numbers}{standard order} $<$. 
\begin{enumerate}
    \item $n < s(n)$ for all $n \in \bbN$
    \item For all $n, m, k \in \bbN$, if $n < m$, then $n + k < m + k$ (equivalently $k + n < k + m$).
    \item For all $n, m, k \in \bbN$, if $n < m$, then $n \times k < m \times k$ (equivalently $k \times n < k \times m$).
    \item the order satisfies trichotomy: for any $n, m \in \bbN$, either $n < m$, $m < n$, or $n = m$.
\end{enumerate}
\end{proposition}

\begin{proof}
    \begin{enumerate}
        \item By \Cref{proposition:natural_numbers_addition_existence}, $n + 1 = s(n)$. Since $1 \in \bbN$, the definition of order gives $n < s(n)$.

        \item If $n < m$, then there exists some $a \in \bbN$ such that 
        $$n + a = m$$
        and hence 
        $$(n+a)+k = m+k.$$
        The \CrefAndHyperref{proposition:natural_numbers_addition_associative_commutative}{associativity and commutativity of addition} imply that the left hand side equals
        $$n+(a+k) = n+(k+a) = (n+k)+a$$
        and hence $n+k < m+k$.

        \item 
        If $n < m$, then there exists some $a \in \bbN$ such that 
        $$n + a = m$$
        and hence 
        $$(n+a) \times k = m \times k.$$
        The \CrefAndHyperref{proposition:natural_numbers_multiplication_associative_commutative_distributive}{distributivity of multiplication} imply that the left hand side equals
        $$n \times k + a \times k$$
        and hence $n \times k < m \times k$.

        \item Since $<$ is irreflexive and asymmetric by \Cref{lemma:standard_order_on_natural_numbers_is_irreflexive} and \Cref{lemma:standard_order_on_natural_numbers_is_asymmetric}, at most one of $n < m$, $m < n$, and $n = m$ can hold for any $n,m \in \bbN$. It thus suffices to show that at least one of these holds.

        Fix $n \in \bbN$. Let $P(m)$ be the statement ``$n = m$, $n < m$, or $m < n$'' for $m \in \bbN$. To show trichotomy is to show that $P(m)$ is true for all $m \in \bbN$. We do so by induction on $m$. In the base case of $m = 1$, either $n = 1$, in which case $n = m$ holds, or $n \neq 1$. We prove the following statement: 
        \begin{center}
            If $n \in \bbN$ is not $1$, then $1 < n$
        \end{center}
        by induction. For $n = 1$, this is vacuously true. Suppose inductively that $1 < k$ holds for some $k \in \bbN$ (and that $1 \neq k$). We show that $1 < s(k)$. We have that $s(k) = k+1$ by \Cref{proposition:natural_numbers_addition_existence}, and the \CrefAndHyperref{proposition:natural_numbers_addition_associative_commutative}{commutativity of addition} equates this with $1+k$, and therefore $1 < s(k)$. Therefore, $P(1)$ holds. If $n = k$, then $n < k+1$. If $n < k$, then there exists some $a \in \bbN$ such that $n + a = k$, so $n +(a+1) = (n+a)+1 = k+1$ and hence $n < k+1$. If 

        Now suppose inductively that $k \in \bbN$ satisfies $P(k)$, i.e. that either $n = k$, $n < k$, or $k < n$. We show that $P(s(k))$ is true, i.e. that $n = s(k)$, $n < s(k)$, or $s(k) < n$. Recall that $s(k) = k+1$ by \Cref{proposition:natural_numbers_addition_existence}. If $n = k$, then $n+1 = k+1$, so $n < k+1$. If $n < k$, then since $k < k+1$ by what we just showed, we have that $n < k+1$ by the \CrefAndHyperrefIfExist{definition:reflexive_symmetric_antisymmetric_transitive_total_binary_relations_on_a_set}{transitivity} of $<$ (\Cref{lemma:standard_order_on_natural_numbers_is_transitive}). If $k < n$, then $k+d = n$ for some $d \in \bbN$. If $d = 1$, then $k+1 = n$, otherwise, $d = s(j)$ for some $j$ by \Cref{lemma:natural_numbers_are_either_1_or_are_successors}. Thus, $k+s(j) = k+(j+1)$. By the \CrefAndHyperrefIfExist{proposition:natural_numbers_addition_associative_commutative}{associativity and commutativity of addition}, this equals $(k+1)+j$, which in turn equals $s(k)+j$ by \Cref{proposition:natural_numbers_addition_existence}. Thus, we have $s(k) + j = n$, so $s(k) < n$. Therefore, $P(m)$ holds by induction and hence $<$ satisfies trichotomy as desired.
    \end{enumerate}
\end{proof}

\begin{proposition} \label{proposition:natural_numbers_well_ordered}
Let $\bbN$ be the set of natural numbers equipped with the \CrefAndHyperrefIfExist{definition:standard_order_on_natural_numbers}{standard order} $<$. Then $(\bbN, <)$ is a \CrefAndHyperrefIfExist{definition:well_ordered_set}{well-ordered set}: every nonempty subset of $\bbN$ has a least element.
\TODO{defien least element}
\end{proposition}

\begin{proof}
Let $S \subseteq \bbN$ be a nonempty subset. Suppose for contradiction that $S$ has no least element. Define
$$L = \{n \in \bbN : \forall m \in S, \ n < m\}.$$
In words, $L$ is the set of natural numbers strictly less than every element of $S$. For $n \in \bbN$, Let $P(n)$ be the statement that $n \in L$. We show that $L = \bbN$ by showing that $P(n)$ holds for all $n \in \bbN$ by induction.

By \Cref{lemma:1_is_least_element_of_natural_numbers}, 
In the case that $n = 1$, suppose for contradiction that $1 \not\in L$. This means that there exists $m \in S$ with $1 \not< m$. By trichotomy on $\bbN$ (\CrefAndHyperrefIfExist{proposition:natural_numbers_order_compatibility}{order compatibility}), either $m = 1$ or $m < 1$. If $m = 1$, then $1 \in S$. Since $1 \leq x$ for all $x \in \bbN$ by \Cref{lemma:1_is_least_element_of_natural_numbers}, this implies that $1$ is the least element of $S$, which is a contradiction. If $m \neq 1$, then $m < 1$ is also impossible by \Cref{lemma:1_is_least_element_of_natural_numbers}. Hence, $1 \in L$, i.e. $P(1)$ holds. 

Suppose inductively that $k \in \bbN$ satisfies $P(k)$. We show that $P(s(k))$ holds, i.e. that $s(k) \in L$. Suppose otherwise for contradiction. In other words, there exists $m \in S$ with $s(k) \not< m$, i.e. $m \leq s(k)$. Suppose for contradiction that $m < s(k)$. Since $k \in L$ by the inductive hypothesis and since $m \in S$, we have $k < m$, so there exists some $a \in \bbN$ such that $m = k+a$. If $a = 1$, then $m = k+1$, contradicting the supposition that $m < s(k) = k+1$ by \Cref{proposition:natural_numbers_addition_existence} and the trichotomy of $<$ (\Cref{proposition:natural_numbers_order_compatibility}). If $a \neq 1$, then $1 < a$ by \Cref{lemma:1_is_least_element_of_natural_numbers} and hence $m < s(k) = k+1 < k+a = m$ by \Cref{proposition:natural_numbers_order_compatibility} and the \CrefAndHyperrefIfExist{proposition:natural_numbers_addition_associative_commutative}{commutativity of addition}, which also violates the trichotomy of $<$. Hence, $m \not< s(k)$ and hence $m = s(k)$. In particular, $s(k) \in S$.

We show that $s(k)$ would have to be the least element of $S$, contradicting the assumption that $S$ has no least element. For any $m' \in S$, since $k \in L$, we have $k < m'$. Let $a \in \bbN$ satisfy $k+a = m'$. If $a = 1$, then $m' = s(k)$ by \Cref{proposition:natural_numbers_addition_existence}. If $a \neq 1$, then $s(k) = k+1 < k+a = m'$ again by \Cref{proposition:natural_numbers_order_compatibility} and the \CrefAndHyperrefIfExist{proposition:natural_numbers_addition_associative_commutative}{commutativity of addition}. In any case, $s(k) \leq m'$, so $s(k)$ is the least element of $S$, yielding the desired contradiction. Hence, $s(k) \in L$, i.e. $P(s(k))$ holds, so $P(n)$ holds for all $n \in \bbN$ by induction. 

In other words, we have shown that $L = \bbN$. Since $S$ is assumed to be nonempty, there exists some $m_0 \in S$. Since $L = \bbN$, we have $m_0 \in L$, meaning $m_0 < m_0$, which contradicts the irreflexivity of $<$\CrefIfExists{lemma:standard_order_on_natural_numbers_is_irreflexive}. Hence, $S$ must in fact have a least element.
\end{proof}


\subsection{Semirings}

\begin{definition} \label{definition:semiring}
    A \hldef{semiring} is a set $R$ equipped with two \CrefAndHyperrefIfExist{definition:n_ary_operation_binary_operation_ternary_operation_on_a_set}{binary operations} $+$ and $\cdot$, called \hldef{addition and multiplication on $R$}, and two distinguished elements \hl{$0_R$} and \hl{$1_R$}, called the \hldef{additive identity of $R$} and the \hldef{multiplicative identity of $R$} respectively, such that
    \begin{enumerate}
        \item $(R, +)$ is a \CrefAndHyperrefIfExist{definition:monoid}{commutative monoid} with identity $0_R$;
        \item $(R, \cdot)$ is a \CrefAndHyperrefIfExist{definition:monoid}{monoid} with identity $1_R$;
        \item multiplication distributes over addition: for all $a, b, c \in R$,
            $$a \cdot (b + c) = a \cdot b + a \cdot c \quad\text{and}\quad (a + b) \cdot c = a \cdot c + b \cdot c;$$
        \item the additive identity annihilates $R$ under multiplication: for all $a \in R$,
            $$0_R \cdot a = 0_R = a \cdot 0_R.$$
    \end{enumerate}
    When the context is clear, we write $0$ and $1$ for $0_R$ and $1_R$, and we write $ab$ for $a \cdot b$.

    We say that the semiring $R$ is \hldef{commutative} if $\cdot$ is a \CrefAndHyperrefIfExist{definition:commutative_binary_operation_on_a_set}{commutative operation}, i.e. that $a \cdot b = b \cdot a$ for all $a,b \in R$.
\end{definition}
\begin{definition} \label{definition:totally_ordered_semiring}
    A \hldef{totally ordered semiring} is an \CrefAndHyperrefIfExist{definition:ordered_semiring}{ordered semiring} whose partial order $\leq$ is a \CrefAndHyperrefIfExist{definition:total_order_linear_order_on_a_set}{total order}.
\end{definition}

\subsection{Natural Numbers with Zero}

\begin{definition} \label{definition:natural_numbers_with_zero}
    The \hldef{set of natural numbers including zero} is the set \hl{$\bbN_0 := \bbN \cup \{0\}$}, where $\bbN$ is the \CrefAndHyperrefIfExist{definition:natural_numbers}{set of natural numbers} and $0$ is a distinguished element not in $\bbN$\CrefIfExists{remark:natural_numbers_zero_inclusive}.

    The binary operation \hldef{addition on $\bbN_0$}, denoted by $+$, is defined by extending \CrefAndHyperrefIfExist{proposition:natural_numbers_addition_existence}{addition on $\bbN$} as follows: for all $n \in \bbN_0$,
    $$0 + n = n + 0 = n,$$
    and for $n, m \in \bbN$, $n + m$ is given by addition in $\bbN$.

    The binary operation \hldef{multiplication on $\bbN_0$}, denoted by $\cdot$, is defined by extending \CrefAndHyperrefIfExist{proposition:natural_numbers_multiplication_existence}{multiplication on $\bbN$} as follows: for all $n \in \bbN_0$,
    $$0 \cdot n = n \cdot 0 = 0,$$
    and for $n, m \in \bbN$, $n \cdot m$ is given by multiplication in $\bbN$.

    The \hldef{standard order on $\bbN_0$} is the binary relation \hl{$\leq$} defined by $0 \leq n$ for all $n \in \bbN_0$, and for $n, m \in \bbN$, $n \leq m$ is given by the \CrefAndHyperrefIfExist{definition:standard_order_on_natural_numbers}{standard order on $\bbN$}. The strict order $<$ is defined as usual by $x < y$ if $x \leq y$ and $x \neq y$, and we write $y > x$ and $y \geq x$ accordingly.
\end{definition}
\begin{proposition} \label{proposition:natural_numbers_with_zero_form_commutative_semiring}
    Let $\bbN_0$ be the \CrefAndHyperrefIfExist{definition:natural_numbers_with_zero}{set of natural numbers including zero} equipped with addition, multiplication, and the standard order. Then:
    \begin{enumerate}
        \item $(\bbN_0, +, \cdot)$ is a \CrefAndHyperrefIfExist{definition:semiring}{commutative semiring};
        \item $(\bbN_0, +)$ is a \CrefAndHyperrefIfExist{definition:cancellative_monoid}{cancellative monoid};
        \item $(\bbN_0, +, \cdot, \leq)$ is a \CrefAndHyperrefIfExist{definition:totally_ordered_semiring}{totally ordered semiring}.
    \end{enumerate}
\end{proposition}

\begin{proof}
\TODO{AI generated, verify}
    We prove each part in turn, using the established properties of \CrefAndHyperrefIfExist{definition:natural_numbers}{$\bbN$}.

    \begin{enumerate}
        \item We verify the semiring axioms.
        \begin{enumerate}
            \item $(\bbN_0, +)$ is a commutative monoid with identity $0$. Associativity and commutativity on $\bbN$ follow from \Cref{proposition:natural_numbers_addition_associative_commutative}, and the cases involving $0$ are immediate from $0 + n = n + 0 = n$. The element $0$ is the additive identity by definition.

            \item $(\bbN_0, \cdot)$ is a monoid with identity $1$. Associativity and commutativity on $\bbN$ follow from \Cref{proposition:natural_numbers_multiplication_associative_commutative_distributive}, and the cases involving $0$ are immediate from $0 \cdot n = n \cdot 0 = 0$. The element $1$ is the multiplicative identity by definition.

            \item Distributivity holds on $\bbN$ by \Cref{proposition:natural_numbers_multiplication_associative_commutative_distributive}, and the cases where any of the three arguments is $0$ are immediate from $0 \cdot x = x \cdot 0 = 0$.

            \item The additive identity annihilates $\bbN_0$ under multiplication: this holds by definition, since $0 \cdot n = n \cdot 0 = 0$ for all $n \in \bbN_0$.
        \end{enumerate}
        Hence $(\bbN_0, +, \cdot)$ is a commutative semiring.

        \item We show that $(\bbN_0, +)$ is cancellative. Let $n, m, k \in \bbN_0$ satisfy $n + k = m + k$.
        \begin{enumerate}
            \item If $k = 0$, then $n = n + 0 = m + 0 = m$.
            \item Suppose $k \in \bbN$. Consider the three possibilities for $n$ and $m$.
            \begin{enumerate}
                \item If $n, m \in \bbN$, then $n + k = m + k$ is an equality in $\bbN$, and \Cref{proposition:natural_numbers_addition_associative_commutative} (cancellativity) gives $n = m$.
                \item If $n = 0$, then $k = m + k$. If $m \in \bbN$, then by the \CrefAndHyperrefIfExist{definition:standard_order_on_natural_numbers}{standard order on $\bbN$} we have $k < m + k$ because $k + m = m + k$\CrefIfExists{proposition:natural_numbers_addition_associative_commutative}. This contradicts $k = m + k$. Hence $m = 0 = n$.
                \item If $m = 0$, the symmetric argument gives $n = 0 = m$.
            \end{enumerate}
        \end{enumerate}
        Thus $n = m$ in all cases, so $(\bbN_0, +)$ is cancellative.

        \item We verify the axioms of a totally ordered semiring.
        \begin{enumerate}
            \item The relation $\leq$ is a total order. Reflexivity, transitivity, antisymmetry, and totality follow from \Cref{proposition:natural_numbers_order_compatibility} for elements of $\bbN$, and the cases involving $0$ are immediate from the definition that $0 \leq n$ for all $n \in \bbN_0$.
            \item Compatibility with addition: if $a \leq b$, then $a + c \leq b + c$ for all $a, b, c \in \bbN_0$. If $a, b, c \in \bbN$, this follows from \Cref{proposition:natural_numbers_order_compatibility}. The remaining cases are straightforward from the definition of $\leq$ on $\bbN_0$.
            \item Compatibility with multiplication: if $a, b \in \bbN_0$ satisfy $a \geq 0$ and $b \geq 0$ (which is true of all elements of $\bbN_0$), then $a b \geq 0$ because the product of any two elements of $\bbN_0$ lies in $\bbN_0$.
        \end{enumerate}
    \end{enumerate}
\end{proof}


% ============================================================
% Section 3: The Integers
% ============================================================
\section{The Integers}

\subsection{The Grothendieck Construction}

\begin{definition} \label{definition:grothendieck_ring_of_a_commutative_semiring}
    Let $R$ be a \CrefAndHyperrefIfExist{definition:semiring}{commutative semiring}. Define a \CrefAndHyperrefIfExist{definition:equivalence_relation_on_a_set}{relation} $\sim$ on $R \times R$ by
    $$(a,b) \sim (c,d) \iff \exists t \in R \text{ such that } a + d + t = b + c + t.$$
    \CrefIfExists{proposition:grothendieck_ring_equivalence_relation}

    The \hldef{Grothendieck ring of $R$} is the \CrefAndHyperrefIfExist{definition:equivalence_class_of_an_element_of_a_set_under_an_equivalence_relation_and_quotient_of_set_by_equivalence_relation}{quotient set}
    $$\hlin{K(R) := (R \times R) / \sim},$$
    whose elements are \CrefAndHyperrefIfExist{definition:equivalence_class_of_an_element_of_a_set_under_an_equivalence_relation_and_quotient_of_set_by_equivalence_relation}{equivalence classes} denoted \hl{$[(a,b)]$}. We define binary operations on $K(R)$, called \hldef{addition and multiplication on $K(R)$}, by
    \begin{align*}
        \hlin{[(a,b)] + [(c,d)]} &\hlin{:= [(a + c,\, b + d)]}, \\
        \hlin{[(a,b)] \cdot [(c,d)]} &\hlin{:= [(ac + bd,\, ad + bc)]},
    \end{align*}
    \footnote{$[(a,b)]$ should be thought of as $a-b$.}
    and distinguished elements
    $$\hlin{0_{K(R)} := [(0_R, 0_R)]}, \qquad \hlin{1_{K(R)} := [(1_R, 0_R)]}.$$
    These operations are well-defined\CrefIfExists{proposition:grothendieck_ring_operations_well_defined} and make $K(R)$ into a \CrefAndHyperrefIfExist{definition:commutative_ring}{commutative unital ring}\CrefIfExists{proposition:grothendieck_ring_is_a_ring}. The additive inverse of $[(a,b)]$ is given by \hl{$-[(a,b)] = [(b,a)]$}.

    There is a natural \hldef{canonical map} \hl{$\iota_R: R \to K(R)$} defined by $\iota_R(r) = [(r, 0_R)]$, which is a homomorphism of semirings. If $R$ is \CrefAndHyperrefIfExist{definition:cancellative_monoid}{cancellative} as an additive monoid, then $\iota_R$ is injective and the defining condition simplifies to $(a,b) \sim (c,d) \iff a + d = b + c$.
    \TODO{write this cancellative equivalence as a lemma}
\end{definition}
\begin{proposition} \label{proposition:grothendieck_ring_equivalence_relation}
Let $R$ be a \CrefAndHyperrefIfExist{definition:semiring}{commutative semiring}. The relation $\sim$ on $R \times R$ defined by $(a,b) \sim (c,d) \iff \exists t \in R \text{ such that } a + d + t = b + c + t$ is an \CrefAndHyperrefIfExist{definition:equivalence_relation_on_a_set}{equivalence relation}.
\end{proposition}

\begin{proof}
\TODO{AI generated, verify}
We verify each axiom of an equivalence relation.

\begin{enumerate}
    \item The relation $\sim$ is reflexive because for any $(a,b) \in R \times R$, we have $a + b + 0_R = b + a + 0_R$ by commutativity of addition in $R$\CrefIfExists{definition:semiring}. Thus $(a,b) \sim (a,b)$.

    \item The relation $\sim$ is symmetric because if $(a,b) \sim (c,d)$, then there exists $t \in R$ such that $a + d + t = b + c + t$. Commutativity of addition then gives $c + b + t = d + a + t$, so $(c,d) \sim (a,b)$.

    \item The relation $\sim$ is transitive because if $(a,b) \sim (c,d)$ and $(c,d) \sim (e,f)$, then there exist $s, t \in R$ such that
    $$a + d + s = b + c + s \quad\text{and}\quad c + f + t = d + e + t.$$
    Adding these two equations gives
    $$a + d + c + f + s + t = b + c + d + e + s + t.$$
    By associativity and commutativity of addition in $R$, this rearranges to
    $$(a + f) + (c + d + s + t) = (b + e) + (c + d + s + t).$$
    Setting $u = c + d + s + t \in R$, we obtain $a + f + u = b + e + u$, so $(a,b) \sim (e,f)$.
\end{enumerate}
\end{proof}
\begin{proposition} \label{proposition:grothendieck_ring_operations_well_defined}
Let $R$ be a \CrefAndHyperrefIfExist{definition:semiring}{commutative semiring} and let $\sim$ be the equivalence relation from \Cref{proposition:grothendieck_ring_equivalence_relation}. If $(a,b) \sim (a',b')$ and $(c,d) \sim (c',d')$ in $R \times R$, then
\begin{enumerate}
    \item $(a + c,\, b + d) \sim (a' + c',\, b' + d')$;
    \item $(ac + bd,\, ad + bc) \sim (a'c' + b'd',\, a'd' + b'c')$.
\end{enumerate}
\end{proposition}

\begin{proof}
\TODO{AI generated, verify}
Because $(a,b) \sim (a',b')$ and $(c,d) \sim (c',d')$, there exist $s, t \in R$ such that
\begin{align}
    a + b' + s &= b + a' + s, \label{eq:witness1} \\
    c + d' + t &= d + c' + t. \label{eq:witness2}
\end{align}

\begin{enumerate}
    \item For addition, add \Cref{eq:witness1} and \Cref{eq:witness2} to obtain
    $$a + b' + c + d' + s + t = b + a' + d + c' + s + t.$$
    By associativity and commutativity of addition in $R$\CrefIfExists{definition:semiring}, this rearranges to
    $$(a + c) + (b' + d') + (s + t) = (b + d) + (a' + c') + (s + t).$$
    Hence $(a + c,\, b + d) \sim (a' + c',\, b' + d')$.

    \item For multiplication, we prove the claim by first showing that the class of the product is unchanged when the first factor is replaced by an equivalent pair, and then that the same holds for the second factor.

    We first replace the first factor. From \Cref{eq:witness1}, multiply both sides by $c$ and by $d$ respectively, using distributivity in $R$\CrefIfExists{definition:semiring}:
    \begin{align}
        ac + b'c + sc &= bc + a'c + sc, \label{eq:m1} \\
        ad + b'd + sd &= bd + a'd + sd. \label{eq:m2}
    \end{align}
    The left-hand side of \Cref{eq:m1} plus the right-hand side of \Cref{eq:m2} equals the right-hand side of \Cref{eq:m1} plus the left-hand side of \Cref{eq:m2}:
    $$(ac + b'c + sc) + (bd + a'd + sd) = (bc + a'c + sc) + (ad + b'd + sd).$$
    By associativity and commutativity of addition, this simplifies to
    $$(ac + bd) + (a'd + b'c) + s(c+d) = (ad + bc) + (a'c + b'd) + s(c+d).$$
    Thus $(ac + bd,\, ad + bc) \sim (a'c + b'd,\, a'd + b'c)$.

    We now replace the second factor. Apply the previous argument with $(c,d)$ in the role of the first factor. Since multiplication is commutative\CrefIfExists{definition:semiring}, we have
    $$(ac + bd,\, ad + bc) = (ca + db,\, cb + da),$$
    and interchanging the roles of $(a,b)$ and $(c,d)$ in Step 1 gives
    $$(ca + db,\, cb + da) \sim (c'a + d'b,\, c'b + d'a).$$
    Reversing the commutativity, $(c'a + d'b,\, c'b + d'a) = (a'c' + b'd',\, a'd' + b'c')$.

    Chaining the two equivalences via transitivity yields
    $$(ac + bd,\, ad + bc) \sim (a'c' + b'd',\, a'd' + b'c'),$$
    as required.
\end{enumerate}
\end{proof}
\begin{proposition} \label{proposition:grothendieck_ring_is_a_ring}
Let $R$ be a \CrefAndHyperrefIfExist{definition:semiring}{commutative semiring} and let $K(R)$ be its \CrefAndHyperrefIfExist{definition:grothendieck_ring_of_a_commutative_semiring}{Grothendieck ring}. Then $K(R)$ is a \CrefAndHyperrefIfExist{definition:commutative_ring}{commutative unital ring}.
\end{proposition}

\begin{proof}
\TODO{AI generated, verify}
Let $\alpha = [(a,b)]$, $\beta = [(c,d)]$, $\gamma = [(e,f)]$ denote arbitrary elements of $K(R)$. We verify each axiom.

\begin{enumerate}
    \item Addition is associative because
    $$(\alpha + \beta) + \gamma = [(a+c+e,\; b+d+f)] = \alpha + (\beta + \gamma),$$
    where the equality uses associativity of addition in $R$\CrefIfExists{definition:semiring}.

    \item Addition is commutative because
    $$\alpha + \beta = [(a+c,\; b+d)] = [(c+a,\; d+b)] = \beta + \alpha,$$
    where the middle equality uses commutativity of addition in $R$\CrefIfExists{definition:semiring}.

    \item The element $0_{K(R)} = [(0_R, 0_R)]$ is the additive identity because for any $\alpha$,
    $$\alpha + 0_{K(R)} = [(a+0_R,\; b+0_R)] = [(a,b)] = \alpha,$$
    where we use that $0_R$ is the additive identity in $R$\CrefIfExists{definition:semiring}. Commutativity of addition then gives $0_{K(R)} + \alpha = \alpha$.

    \item Every element $\alpha$ has an additive inverse $-\alpha := [(b,a)]$ because
    $$\alpha + (-\alpha) = [(a+b,\; b+a)].$$
    Since $a+b = b+a$ by commutativity of addition in $R$, we have $(a+b,\; b+a) \sim (0_R, 0_R)$ with witness $t = 0_R$:
    $$(a+b) + 0_R + 0_R = (b+a) + 0_R + 0_R.$$
    Hence $\alpha + (-\alpha) = 0_{K(R)}$. By commutativity of addition, $(-\alpha) + \alpha = 0_{K(R)}$.

    \item Multiplication is associative. Compute
    \begin{align*}
    (\alpha \beta) \gamma &= [(ac+bd,\; ad+bc)] \cdot [(e,f)] \\
    &= [((ac+bd)e + (ad+bc)f,\; (ac+bd)f + (ad+bc)e)].
    \end{align*}
    Using distributivity, associativity, and commutativity of multiplication in $R$\CrefIfExists{definition:semiring}, this equals
    $$[(ace + adf + bcf + bde,\; acf + ade + bce + bdf)].$$
    On the other hand,
    \begin{align*}
    \alpha (\beta \gamma) &= [(a,b)] \cdot [(ce + df,\; cf + de)] \\
    &= [(a(ce+df) + b(cf+de),\; a(cf+de) + b(ce+df))] \\
    &= [(ace + adf + bcf + bde,\; acf + ade + bce + bdf)],
    \end{align*}
    where the expansion again uses distributivity and the properties of $R$. Since both expressions are represented by the same pair, they are equal in $K(R)$.

    \item Multiplication is commutative because
    \begin{align*}
    \alpha \beta &= [(ac+bd,\; ad+bc)] \\
    &= [(ca+db,\; da+cb)] \quad\text{(commutativity of $\cdot$ in $R$)}\\
    &= [(ca+db,\; cb+da)] \quad\text{(commutativity of $+$ in $R$)}\\
    &= \beta \alpha.
    \end{align*}

    \item The element $1_{K(R)} = [(1_R, 0_R)]$ is the multiplicative identity because
    $$\alpha \cdot 1_{K(R)} = [(a \cdot 1_R + b \cdot 0_R,\; a \cdot 0_R + b \cdot 1_R)] = [(a,\; b)] = \alpha,$$
    where we use that $1_R$ is the multiplicative identity in $R$ and that $0_R$ annihilates under multiplication\CrefIfExists{definition:semiring}. Commutativity of multiplication then gives $1_{K(R)} \cdot \alpha = \alpha$.

    \item Multiplication distributes over addition. We verify left distributivity; right distributivity follows from commutativity.
    \begin{align*}
    \alpha (\beta + \gamma) &= [(a,b)] \cdot [(c+e,\; d+f)] \\
    &= [(a(c+e) + b(d+f),\; a(d+f) + b(c+e))] \\
    &= [(ac+ae+bd+bf,\; ad+af+bc+be)].
    \end{align*}
    On the other hand,
    \begin{align*}
    \alpha\beta + \alpha\gamma &= [(ac+bd,\; ad+bc)] + [(ae+bf,\; af+be)] \\
    &= [(ac+bd+ae+bf,\; ad+bc+af+be)].
    \end{align*}
    The two expressions are represented by the same pair (by commutativity of addition in $R$), so $\alpha(\beta+\gamma) = \alpha\beta + \alpha\gamma$.
\end{enumerate}
Thus $K(R)$ satisfies all axioms of a commutative unital ring.
\end{proof}
\begin{theorem} \label{theorem:grothendieck_ring_universal_property}
\TODO{AI generated, verify}
\TODO{semiring homomorphism}
Let $R$ be a \CrefAndHyperrefIfExist{definition:semiring}{commutative semiring} with \CrefAndHyperrefIfExist{definition:grothendieck_ring_of_a_commutative_semiring}{Grothendieck ring} $K(R)$ and canonical map $\iota_R: R \to K(R)$. For every \CrefAndHyperrefIfExist{definition:commutative_ring}{commutative unital ring} $S$ and every semiring homomorphism $\phi: R \to S$, there exists a unique \CrefAndHyperrefIfExist{definition:ring_homomorphism}{ring homomorphism} $\widetilde{\phi}: K(R) \to S$ such that $\widetilde{\phi} \circ \iota_R = \phi$.
\end{theorem}

\begin{proof}
\TODO{AI generated, verify}
We define $\widetilde{\phi}: K(R) \to S$ by $\widetilde{\phi}([(a,b)]) = \phi(a) - \phi(b)$, where the subtraction on the right is taken in the ring $S$.

We first check that $\widetilde{\phi}$ is well-defined. Suppose $(a,b) \sim (c,d)$, so there exists $t \in R$ such that $a + d + t = b + c + t$. Applying $\phi$ to both sides and using that $\phi$ respects addition, we obtain
$$\phi(a) + \phi(d) + \phi(t) = \phi(b) + \phi(c) + \phi(t)$$
in $S$. Subtracting $\phi(b) + \phi(d) + \phi(t)$ from both sides gives $\phi(a) - \phi(b) = \phi(c) - \phi(d)$. Hence $\widetilde{\phi}([(a,b)]) = \widetilde{\phi}([(c,d)])$.

We now verify that $\widetilde{\phi}$ is a ring homomorphism. Let $\alpha = [(a,b)], \beta = [(c,d)] \in K(R)$.

\begin{enumerate}
    \item The map $\widetilde{\phi}$ preserves addition. We have $\alpha + \beta = [(a+c,\, b+d)]$, so
    $$\widetilde{\phi}(\alpha + \beta) = \phi(a+c) - \phi(b+d).$$
    Since $\phi$ respects addition, this equals $(\phi(a) + \phi(c)) - (\phi(b) + \phi(d))$. In the ring $S$, rearranging yields $(\phi(a) - \phi(b)) + (\phi(c) - \phi(d)) = \widetilde{\phi}(\alpha) + \widetilde{\phi}(\beta)$.

    \item The map $\widetilde{\phi}$ preserves multiplication. We have $\alpha \beta = [(ac + bd,\, ad + bc)]$, so
    $$\widetilde{\phi}(\alpha \beta) = \phi(ac + bd) - \phi(ad + bc).$$
    Since $\phi$ respects addition and multiplication, this equals $(\phi(a)\phi(c) + \phi(b)\phi(d)) - (\phi(a)\phi(d) + \phi(b)\phi(c))$. Factoring in the ring $S$, we compute
    \begin{align*}
    \widetilde{\phi}(\alpha)\widetilde{\phi}(\beta) &= (\phi(a) - \phi(b))(\phi(c) - \phi(d)) \\
    &= \phi(a)\phi(c) - \phi(a)\phi(d) - \phi(b)\phi(c) + \phi(b)\phi(d) \\
    &= (\phi(a)\phi(c) + \phi(b)\phi(d)) - (\phi(a)\phi(d) + \phi(b)\phi(c)) \\
    &= \widetilde{\phi}(\alpha \beta).
    \end{align*}
    Here the minus signs distribute using the ring axioms of $S$\CrefIfExists{definition:commutative_ring}.

    \item The map $\widetilde{\phi}$ preserves unity. We have $\widetilde{\phi}(1_{K(R)}) = \widetilde{\phi}([(1_R, 0_R)]) = \phi(1_R) - \phi(0_R)$. Since $\phi(0_R) = 0_S$ and $\phi(1_R) = 1_S$ for a semiring homomorphism, we obtain $\widetilde{\phi}(1_{K(R)}) = 1_S - 0_S = 1_S$.
\end{enumerate}

Thus $\widetilde{\phi}$ is a ring homomorphism. Moreover, for any $r \in R$,
$$\widetilde{\phi}(\iota_R(r)) = \widetilde{\phi}([(r, 0_R)]) = \phi(r) - \phi(0_R) = \phi(r) - 0_S = \phi(r),$$
so $\widetilde{\phi} \circ \iota_R = \phi$.

For uniqueness, suppose $\psi: K(R) \to S$ is any ring homomorphism with $\psi \circ \iota_R = \phi$. For any $[(a,b)] \in K(R)$, write $[(a,b)] = [(a,0_R)] + [(0_R,b)] = \iota_R(a) + (-\iota_R(b))$. Since $\psi$ respects addition and additive inverses,
$$\psi([(a,b)]) = \psi(\iota_R(a)) + \psi(-\iota_R(b)) = \phi(a) - \phi(b) = \widetilde{\phi}([(a,b)]).$$
Hence $\psi = \widetilde{\phi}$, establishing uniqueness.
\end{proof}

\subsection{The Ring of Integers}

\begin{definition} \label{definition:ring_of_integers_from_natural_numbers}
Let $\bbN$ be the set of natural numbers from \CrefAndHyperrefIfExist{definition:natural_numbers}{natural numbers}. Define an \CrefAndHyperrefIfExist{definition:equivalence_relation_on_a_set}{equivalence relation} $\sim$ on $\bbN \times \bbN$ by:
$$(a, b) \sim (c, d) \iff a + d = b + c.$$
\CrefIfExists{proposition:integer_equivalence_relation}
The \hldef{set of integers} is the \CrefAndHyperrefIfExist{definition:equivalence_class_of_an_element_of_a_set_under_an_equivalence_relation_and_quotient_of_set_by_equivalence_relation}{quotient set} \hl{$\bbZ := (\bbN \times \bbN) / \sim$}. The \CrefAndHyperrefIfExist{definition:equivalence_class_of_an_element_of_a_set_under_an_equivalence_relation_and_quotient_of_set_by_equivalence_relation}{equivalence class} of $(a, b)$ is denoted \hl{$[(a, b)]$} and intuitively represents the integer $a - b$.

There are binary operators on $\bbZ$, called \hldef{addition and multiplication on $\bbZ$} defined by:
$$\hlin{[(a, b)] + [(c, d)] := [(a + c, b + d)],}$$
$$\hlin{[(a, b)] \times [(c, d)] := [(a \times c + b \times d, b \times c + a \times d)].}$$
\TODO{define $2$}
There is also a unary operator on $\bbZ$, called \hldef{additive inverse on $\bbZ$} defined by:
$$\hlin{-[(a, b)] := [(b, a)].}$$
Equipped with these operators and the elements \hl{$0_\bbZ := [(1, 1)]$} and \hl{$1_\bbZ := [(2, 1)]$}, $\bbZ$ is a \CrefAndHyperrefIfExist{definition:commutative_ring}{commutative unital ring}\CrefIfExists{proposition:integers_form_ring} and in fact forms an \CrefAndHyperrefIfExist{definition:ordered_ring}{ordered ring}\CrefIfExists{proposition:integers_form_ordered_ring}.
\end{definition}

\begin{proposition} \label{proposition:integer_equivalence_relation}
Let $\bbN$ be the \CrefAndHyperrefIfExist{definition:natural_numbers}{set of natural numbers}. The relation $\sim$ on $\bbN \times \bbN$ defined by $(a, b) \sim (c, d) \iff a + d = b + c$ is an \CrefAndHyperrefIfExist{definition:equivalence_relation_on_a_set}{equivalence relation}.
\end{proposition}

\begin{proof}
For any $a,b \in \bbN$, $(a, b) \sim (a, b)$ because $a + b = b + a$ by \CrefAndHyperrefIfExist{proposition:natural_numbers_addition_associative_commutative}{commutativity of addition on $\bbN$}. If $(a, b) \sim (c, d)$ for $a,b,c,d \in \bbN$, then $a + d = b + c$, so $c + b = d + a$, hence $(c, d) \sim (a, b)$. If $(a, b) \sim (c, d)$ and $(c, d) \sim (e, f)$ for $a,b,c,d,e,f \in \bbN$, then $a + d = b + c$ and $c + f = d + e$. Adding these equations gives $(a + d) + (c + f) = (b + c) + (d + e)$. By \CrefAndHyperrefIfExist{proposition:natural_numbers_addition_associative_commutative}{associativity and commutative of addition}, one eventually shows that $(a+f)+(c+d) = (b+e)+(c+d)$, so $(a, b) \sim (e, f)$.
\end{proof}

\begin{proposition} \label{proposition:integer_operations_well_defined}
Let $\bbZ$ be the set of integers as constructed in \Cref{definition:integers_from_natural_numbers}. The addition and multiplication operations on $\bbZ$ are well-defined with respect to the \CrefAndHyperrefIfExist{definition:equivalence_relation_on_a_set}{equivalence relation} $\sim$. That is, if $(a, b) \sim (a', b')$ and $(c, d) \sim (c', d')$, then:
\begin{enumerate}
    \item $[(a + c, b + d)] = [(a' + c', b' + d')]$,
    \item $[(a \times c + b \times d, b \times c + a \times d)] = [(a' \times c' + b' \times d', b' \times c' + a' \times d')]$.
\end{enumerate}
\end{proposition}

\begin{proof}
    \TODO{AI generated, verify}
From $(a, b) \sim (a', b')$ we have $a + b' = b + a'$ in $\bbN$, and from $(c, d) \sim (c', d')$ we have $c + d' = d + c'$ in $\bbN$. We write these as:
\begin{align*}
\text{(E1)}:\quad a + b' = b + a' \\
\text{(E2)}:\quad c + d' = d + c'
\end{align*}

We first show that addition on $\bbZ$ is well-defined. We need to show that $(a+c) + (b'+d') = (b+d) + (a'+c')$ in $\bbN$. Adding the equalities (E1) and (E2) gives
$$(a + b') + (c + d') = (b + a') + (d + c').$$
By \Cref{proposition:natural_numbers_addition_associative_commutative}, addition on $\bbN$ is associative and commutative, so the left side rearranges to $(a+c) + (b'+d')$ and the right side to $(b+d) + (a'+c')$. Hence $[(a+c, b+d)] = [(a'+c', b'+d')]$.

Next we show that multiplication on $\bbZ$ is well-defined. We need to show that 
$$(a \times c + b \times d) + (b' \times c' + a' \times d') = (b \times c + a \times d) + (a' \times c' + b' \times d')$$
in $\bbN$. Let $X = a \times d' + b \times c' + a' \times c + b' \times d$. We add $X$ to both sides and verify equality.

We compute the left side plus $X$:
\begin{align*}
&\quad (a \times c + b \times d + b' \times c' + a' \times d') + X \\
&= a \times c + b \times d + b' \times c' + a' \times d' + a \times d' + b \times c' + a' \times c + b' \times d \\
&= (a \times c + a \times d') + (b \times d + b \times c') + (b' \times c' + b' \times d) + (a' \times d' + a' \times c) && \text{[\Cref{proposition:natural_numbers_addition_associative_commutative}]} \\
&= a \times (c + d') + b \times (d + c') + b' \times (c' + d) + a' \times (d' + c) && \text{[\Cref{proposition:natural_numbers_multiplication_associative_commutative_distributive}]} \\
&= a \times (d + c') + b \times (d + c') + b' \times (d + c') + a' \times (d + c') && \text{[E2, comm of $+$]} \\
&= (a + b + b' + a') \times (d + c') && \text{[\Cref{proposition:natural_numbers_multiplication_associative_commutative_distributive}]} \\
&= ((a + b') + (b + a')) \times (d + c') && \text{[\Cref{proposition:natural_numbers_addition_associative_commutative}]} \\
&= ((b + a') + (b + a')) \times (d + c') && \text{[E1: $a+b'=b+a'$]} \\
&= (2 \times (b + a')) \times (d + c'). && \text{[$2 \times x = x + x$ for $x \in \bbN$]}
\end{align*}

We compute the right side plus $X$:
\begin{align*}
&\quad (b \times c + a \times d + a' \times c' + b' \times d') + X \\
&= b \times c + a \times d + a' \times c' + b' \times d' + a \times d' + b \times c' + a' \times c + b' \times d \\
&= (b \times c + b \times c') + (a \times d + a \times d') + (a' \times c' + a' \times c) + (b' \times d' + b' \times d) && \text{[\Cref{proposition:natural_numbers_addition_associative_commutative}]} \\
&= b \times (c + c') + a \times (d + d') + a' \times (c' + c) + b' \times (d' + d) && \text{[\Cref{proposition:natural_numbers_multiplication_associative_commutative_distributive}]} \\
&= b \times (c + c') + (b + a') \times (d + d') + a' \times (c + c') + b' \times (d + d') && \text{[E1, comm of $+$]} \\
&= (b + a') \times (c + c') + (b + a') \times (d + d') && \text{[\Cref{proposition:natural_numbers_multiplication_associative_commutative_distributive}]} \\
&= (b + a') \times (c + c' + d + d') && \text{[\Cref{proposition:natural_numbers_multiplication_associative_commutative_distributive}]} \\
&= (b + a') \times ((c + d') + (c' + d)) && \text{[\Cref{proposition:natural_numbers_addition_associative_commutative}]} \\
&= (b + a') \times ((d + c') + (d + c')) && \text{[E2, comm of $+$]} \\
&= (b + a') \times (2 \times (d + c')) \\
&= (2 \times (b + a')) \times (d + c'). && \text{[\Cref{proposition:natural_numbers_multiplication_associative_commutative_distributive}]}
\end{align*}

Both sides plus $X$ equal $(2 \times (b + a')) \times (d + c')$. By \Cref{proposition:natural_numbers_addition_associative_commutative}, addition on $\bbN$ is cancellative (if $u + X = v + X$, then $u = v$), so the target equality holds. Hence multiplication is well-defined.
\end{proof}

\begin{proposition} \label{proposition:integers_form_ring}
Let $\bbZ$ be the set of integers with addition and multiplication from \CrefAndHyperrefIfExist{definition:ring_of_integers_from_natural_numbers}{integers from natural numbers}. Then $\bbZ$ is a \CrefAndHyperrefIfExist{definition:commutative_ring}{commutative ring with unity}. Specifically:
\begin{enumerate}
    \item $(\bbZ, +)$ is an \CrefAndHyperrefIfExist{definition:group}{abelian group} with identity $0_\bbZ = [(1, 1)]$.
    \item $(\bbZ, \times)$ is a \CrefAndHyperrefIfExist{definition:monoid}{commutative monoid} with identity $1_\bbZ = [(2, 1)]$.
    \item Multiplication distributes over addition.
    \item Every element has an additive inverse: $-[(a, b)] = [(b, a)]$.
\end{enumerate}
Moreover, $\bbZ$ is an \CrefAndHyperrefIfExist{definition:domain_integral_domain}{integral domain}, i.e. it has no \CrefAndHyperrefIfExist{definition:zero_divisor_of_a_ring}{zero divisors}.
\end{proposition}

\begin{proof}
Let $\alpha = [(a,b)]$, $\beta = [(c,d)]$, $\gamma = [(e,f)]$ denote arbitrary elements of $\bbZ$. Recall various properties of addition and multiplication on $\bbN$ from \Cref{proposition:natural_numbers_addition_associative_commutative} and \Cref{proposition:natural_numbers_multiplication_associative_commutative_distributive}. 

\begin{enumerate}
    \item We show that $(\bbZ, +)$ is an abelian group.

    \begin{itemize}
        
        \item We show that $+$ is associative.  We have $\alpha + \beta = [(a+c, b+d)]$, so
        $$(\alpha + \beta) + \gamma = [(a+c+e, b+d+f)].$$
        Similarly, $\beta + \gamma = [(c+e, d+f)]$, so
        $$\alpha + (\beta + \gamma) = [(a+(c+e), b+(d+f))].$$
        By associativity of addition on $\bbN$, $a+c+e = a+(c+e)$ and $b+d+f = b+(d+f)$, hence the pairs are identical and $(\alpha+\beta)+\gamma = \alpha+(\beta+\gamma)$.

        \item We show that $+$ is commutative. We have $\alpha + \beta = [(a+c, b+d)] = [(c+a, d+b)] = \beta + \alpha$, where the middle equality uses commutativity of addition on $\bbN$.

        \item We show that $0_\bbZ = [(1,1)]$ is the additive identity for $+$. For any $\alpha = [(a,b)]$,
        $$\alpha + 0_\bbZ = [(a+1, b+1)].$$
        To see that $[(a+1, b+1)] = [(a,b)]$, we check the equivalence condition: $(a+1) + b = a + (b+1)$, which holds by associativity of addition on $\bbN$. By commutativity of addition on $\bbZ$ (proven above), $0_\bbZ + \alpha = \alpha$ as well.

        \item We show that $-\alpha = [(b,a)]$ is the additive inverse of $\alpha$. Indeed,
        $$\alpha + (-\alpha) = [(a+b, b+a)].$$
        To see that $[(a+b, b+a)] = [(1,1)] = 0_\bbZ$, we check: $(a+b) + 1 = (b+a) + 1$, which holds by commutativity of addition on $\bbN$ ($a+b = b+a$). By commutativity of addition on $\bbZ$, $(-\alpha) + \alpha = 0_\bbZ$ as well.

    \end{itemize}


    \item We show that $(\bbZ, \times)$ is a commutative monoid.

    \begin{itemize}
    
        \item We show that $\times$ is commutative. Compute
        $$\alpha \times \beta = [(ac + bd, bc + ad)],$$
        $$\beta \times \alpha = [(ca + db, da + cb)].$$
        To see these are equal, check: $(ac+bd) + (da+cb) = (bc+ad) + (ca+db)$. The left side rearranges to $ac+bd+ad+bc$ and the right side to $bc+ad+ac+bd$, which are equal by associativity and commutativity of addition on $\bbN$.

        \item We show that $1_\bbZ = [(2,1)]$ is the multiplicative identity. Compute
        $$\alpha \times 1_\bbZ = [(a \times 2 + b \times 1, b \times 2 + a \times 1)] = [(2a + b, 2b + a)].$$
        To see that $[(2a+b, 2b+a)] = [(a,b)]$, we check the equivalence condition:
        $$(2a + b) + b = 2a + 2b = a + (2b + a),$$
        where $2a + 2b = a + a + b + b$ and $a + (2b + a) = a + b + b + a$, which are equal by associativity and commutativity of addition on $\bbN$. That $1_\bbZ \times \alpha = \alpha$ follows from the commutativity of $\times$.

        \item We show that $\times$ is associative. We compute $(\alpha \times \beta) \times \gamma$. First,
        $$\alpha \times \beta = [(ac + bd, bc + ad)].$$
        Then
        \begin{align*}
        (\alpha \times \beta) \times \gamma &= [((ac+bd)e + (bc+ad)f,\; (ac+bd)f + (bc+ad)e)] \\
        &= [(ace + bde + bcf + adf,\; acf + bdf + bce + ade)],
        \end{align*}
        where the expansion uses distributivity on $\bbN$.

        Next, we compute $\alpha \times (\beta \times \gamma)$. First,
        $$\beta \times \gamma = [(ce + df, cf + de)].$$
        Then
        \begin{align*}
        \alpha \times (\beta \times \gamma) &= [(a(ce+df) + b(cf+de),\; a(cf+de) + b(ce+df))] \\
        &= [(ace + adf + bcf + bde,\; acf + ade + bce + bdf)],
        \end{align*}
        again using distributivity on $\bbN$.

        To check equivalence, we need the sum of the first component of $(\alpha\times\beta)\times\gamma$ and the second component of $\alpha\times(\beta\times\gamma)$ to equal the sum of the second component of $(\alpha\times\beta)\times\gamma$ and the first component of $\alpha\times(\beta\times\gamma)$:
        \begin{align*}
        &\quad (ace + bde + bcf + adf) + (acf + ade + bce + bdf) \\
        &= ace + bde + bcf + adf + acf + ade + bce + bdf \\
        &= acf + bdf + bce + ade + ace + adf + bcf + bde \\
        &= (acf + bdf + bce + ade) + (ace + adf + bcf + bde),
        \end{align*}
        where the rearrangement uses only associativity and commutativity of addition on $\bbN$. Hence $(\alpha\times\beta)\times\gamma = \alpha\times(\beta\times\gamma)$.

    \end{itemize}

    \item We show that multiplication is distributive over addition, i.e. that $\alpha \times (\beta + \gamma) = \alpha \times \beta + \alpha \times \gamma$.
    Compute $\beta + \gamma = [(c+e, d+f)]$, so
    $$\alpha \times (\beta + \gamma) = [(a(c+e) + b(d+f),\; a(d+f) + b(c+f))].$$
    On the other hand,
    $$\alpha \times \beta = [(ac+bd, bc+ad)], \quad \alpha \times \gamma = [(ae+bf, af+be)],$$
    so
    $$\alpha \times \beta + \alpha \times \gamma = [(ac+bd+ae+bf,\; bc+ad+af+be)].$$
    To check equivalence, we need:
    $$a(c+e) + b(d+f) + (bc+ad+af+be) = (a(d+f) + b(c+f)) + (ac+bd+ae+bf).$$
    Expanding using distributivity on $\bbN$: the left side becomes $ac+ae+bd+bf+bc+ad+af+be$, and the right side becomes $ad+af+bc+bf+ac+bd+ae+be$. These are equal by associativity and commutativity of addition on $\bbN$.

    \item We show that $\bbZ$ has no zero divisors. Suppose that $\alpha \times \beta = 0_\bbZ = [(1,1)]$. This means
    $$(ac + bd, bc + ad) \sim (1, 1),$$
    i.e., $ac + bd + 1 = bc + ad + 1$ in $\bbN$, which simplifies to $ac + bd = bc + ad$ by \CrefAndHyperrefIfExist{proposition:natural_numbers_addition_associative_commutative}{cancellativity of addition on $\bbN$}. We consider cases using \CrefAndHyperrefIfExist{proposition:natural_numbers_order_compatibility}{trichotomy of order on $\bbN$}:
    \begin{enumerate}
        \item If $a = b$, then $\alpha = [(a,a)] = [(1,1)] = 0_\bbZ$, so we are done.
        \item If $c = d$, then $\beta = [(c,c)] = [(1,1)] = 0_\bbZ$, so we are done.
        \item If $a > b$, then by the definition of the \CrefAndHyperrefIfExist{definition:standard_order_on_natural_numbers}{standard order} on $\bbN$ there exists $x \in \bbN$ such that $a = b + x$. Substituting into $ac + bd = bc + ad$:
        $$(b+x)c + bd = bc + (b+x)d.$$
        Expanding by distributivity on $\bbN$: $bc + xc + bd = bc + bd + xd$. By cancellativity of addition on $\bbN$ (canceling $bc + bd$ from both sides), $xc = xd$. Since $x \in \bbN$ and $x \geq 1$, by \CrefAndHyperrefIfExist{proposition:cancellativity_of_multiplication_in_the_natural_numbers}{cancellativity of multiplication on $\bbN$}, $c = d$. Hence $\beta = [(c,c)] = 0_\bbZ$.
        \item If $b > a$, then there exists $x \in \bbN$ with $b = a + x$. Substituting into $ac + bd = bc + ad$ yields
        $$ac + (a+x)d = (a+x)c + ad.$$
        Expanding, we have $ac + ad + xd = ac + xc + ad$. Canceling $ac + ad$ from both sides, we have $xd = xc$. Since $x \geq 1$, cancellativity of multiplication gives $d = c$. Hence $\beta = [(c,c)] = 0_\bbZ$.
    \end{enumerate}
    In all cases, either $\alpha = 0_\bbZ$ or $\beta = 0_\bbZ$.
\end{enumerate}
\end{proof}


\subsection{Universal Properties of \texorpdfstring{$\mathbb Z$}{Z}}

\begin{proposition} \label{proposition:grothendieck_ring_of_natural_numbers_with_zero}
\TODO{AI generated, verify}
    Let $\bbN_0$ be the \CrefAndHyperrefIfExist{definition:natural_numbers_with_zero}{set of natural numbers including zero}, regarded as a \CrefAndHyperrefIfExist{definition:semiring}{semiring}\CrefIfExists{proposition:natural_numbers_with_zero_form_commutative_semiring}. Let $K(\bbN_0)$ be the \CrefAndHyperrefIfExist{definition:grothendieck_ring_of_a_commutative_semiring}{Grothendieck ring} of $\bbN_0$, and let $\bbZ$ be the \CrefAndHyperrefIfExist{definition:ring_of_integers_from_natural_numbers}{ring of integers} as constructed from $\bbN$. Then there is an isomorphism of rings
    $$\Phi: K(\bbN_0) \xrightarrow{\cong} \bbZ.$$
    Consequently, $\bbZ$ is (up to isomorphism) the Grothendieck ring of $\bbN_0$.
\end{proposition}

\begin{proof}
\TODO{AI generated, verify}
    We construct an explicit pair of inverse ring homomorphisms.

    \textbf{Construction of $\Phi$.} Define a map $F: \bbN_0 \times \bbN_0 \to \bbZ$ by
    $$F(a,b) = [(a+1,\, b+1)],$$
    where $a+1, b+1 \in \bbN$ are the successors of $a$ and $b$ respectively (so $0+1 = 1$, and for $n \in \bbN$, $n+1$ is the usual successor).

    We first verify that $F$ respects the Grothendieck equivalence relation on $\bbN_0 \times \bbN_0$. Suppose $(a,b) \sim (c,d)$ in $\bbN_0 \times \bbN_0$, i.e. there exists $t \in \bbN_0$ such that $a + d + t = b + c + t$. Adding $2$ to both sides gives
    $$(a+1) + (d+1) + t = (b+1) + (c+1) + t.$$
    Since $(\bbN_0,+)$ is \CrefAndHyperrefIfExist{definition:cancellative_monoid}{cancellative}\CrefIfExists{proposition:natural_numbers_with_zero_form_commutative_semiring}, the witness $t$ can be cancelled, yielding
    $$(a+1) + (d+1) = (b+1) + (c+1).$$
    Thus $(a+1,b+1) \sim (c+1,d+1)$ in $\bbN \times \bbN$, so $F(a,b) = F(c,d)$. Hence $F$ descends to a well-defined map $\Phi: K(\bbN_0) \to \bbZ$.

    \textbf{$\Phi$ is a ring homomorphism.} Let $[(a,b)], [(c,d)] \in K(\bbN_0)$.
    \begin{enumerate}
        \item Addition. By definition of addition in $K(\bbN_0)$,
        $$\Phi([(a,b)] + [(c,d)]) = \Phi([(a+c, b+d)]) = [(a+c+1,\, b+d+1)].$$
        On the other hand,
        $$\Phi([(a,b)]) + \Phi([(c,d)]) = [(a+1,b+1)] + [(c+1,d+1)] = [(a+c+2,\, b+d+2)].$$
        These two classes are equal in $\bbZ$ because
        $$(a+c+2) + (b+d+1) = a+b+c+d+3 = (a+c+1) + (b+d+2)$$
        by associativity and commutativity of addition in $\bbN$\CrefIfExists{proposition:natural_numbers_addition_associative_commutative}. Hence $\Phi$ preserves addition.

        \item Multiplication. By definition of multiplication in $K(\bbN_0)$,
        $$\Phi([(a,b)] \cdot [(c,d)]) = \Phi([(ac+bd,\, ad+bc)]) = [(ac+bd+1,\, ad+bc+1)].$$
        On the other hand,
        \begin{align*}
        \Phi([(a,b)]) \cdot \Phi([(c,d)]) &= [(a+1,b+1)] \cdot [(c+1,d+1)] \\
        &= [((a+1)(c+1)+(b+1)(d+1),\; (a+1)(d+1)+(b+1)(c+1))].
        \end{align*}
        Expanding by distributivity in $\bbN$\CrefIfExists{proposition:natural_numbers_multiplication_associative_commutative_distributive}:
        \begin{align*}
        (a+1)(c+1)+(b+1)(d+1) &= ac+a+c+1+bd+b+d+1 = ac+bd + a+b+c+d+2, \\
        (a+1)(d+1)+(b+1)(c+1) &= ad+a+d+1+bc+b+c+1 = ad+bc + a+b+c+d+2.
        \end{align*}
        Now compare $(ac+bd+1,\, ad+bc+1)$ with $(ac+bd+a+b+c+d+2,\; ad+bc+a+b+c+d+2)$. We have
        $$(ac+bd+1) + (ad+bc+a+b+c+d+2) = ac+bd+ad+bc + a+b+c+d+3$$
        and
        $$(ad+bc+1) + (ac+bd+a+b+c+d+2) = ad+bc+ac+bd + a+b+c+d+3,$$
        which are equal by commutativity of addition in $\bbN$. Hence the two classes coincide, so $\Phi$ preserves multiplication.

        \item Unity. The multiplicative identity of $K(\bbN_0)$ is $[(1,0)]$. We have $\Phi([(1,0)]) = [(2,1)] = 1_\bbZ$.
    \end{enumerate}
    Thus $\Phi$ is a ring homomorphism.

    \textbf{Construction of the inverse $\Psi$.} Define $\Psi: \bbZ \to K(\bbN_0)$ by
    $$\Psi([(a,b)]) = [(a,b)],$$
    where $(a,b) \in \bbN \times \bbN$ is regarded as a pair in $\bbN_0 \times \bbN_0$ via the inclusion $\bbN \subset \bbN_0$. This is well-defined because if $(a,b) \sim (c,d)$ in $\bbN \times \bbN$ (i.e. $a+d = b+c$), then the same equality holds in $\bbN_0$, so $(a,b) \sim (c,d)$ in $\bbN_0 \times \bbN_0$. Moreover, $\Psi$ is a ring homomorphism because the ring operations in $\bbZ$ and $K(\bbN_0)$ are defined by the same formulas on representatives.

    \textbf{$\Phi$ and $\Psi$ are inverses.}
    \begin{enumerate}
        \item $\Phi \circ \Psi = \id_\bbZ$. For $[(a,b)] \in \bbZ$,
        $$\Phi(\Psi([(a,b)])) = \Phi([(a,b)]) = [(a+1,b+1)].$$
        In $\bbZ$, we have $(a+1,b+1) \sim (a,b)$ because
        $$(a+1) + b = a + b + 1 = (b+1) + a$$
        by commutativity of addition in $\bbN$. Thus $[(a+1,b+1)] = [(a,b)]$, so $\Phi \circ \Psi = \id_\bbZ$.

        \item $\Psi \circ \Phi = \id_{K(\bbN_0)}$. For $[(a,b)] \in K(\bbN_0)$,
        $$\Psi(\Phi([(a,b)])) = \Psi([(a+1,b+1)]) = [(a+1,b+1)].$$
        We claim that $(a+1,b+1) \sim (a,b)$ in $K(\bbN_0)$. Indeed, by commutativity of addition in $\bbN_0$,
        $$(a+1) + b = a + 1 + b = a + (b+1),$$
        so with $t = 0$ we have $(a+1)+b+0 = a+(b+1)+0$. Hence $[(a+1,b+1)] = [(a,b)]$, and $\Psi \circ \Phi = \id_{K(\bbN_0)}$.
    \end{enumerate}

    Therefore $\Phi$ is an isomorphism of rings with inverse $\Psi$.
\end{proof}
\begin{proposition} \label{proposition:positive_integers_correspond_to_natural_numbers}
\TODO{AI generated, verify}
    Let $\bbZ$ be the \CrefAndHyperrefIfExist{definition:integers_from_natural_numbers}{ring of integers} and let $\bbZ_{> 0} \subset \bbZ$ be its \CrefAndHyperrefIfExist{definition:positive_cone_in_an_ordered_monoid}{positive cone}. Let $\bbN$ be the \CrefAndHyperrefIfExist{definition:natural_numbers}{set of natural numbers}, regarded as a \CrefAndHyperrefIfExist{definition:totally_ordered_semiring}{totally ordered semiring} via the \CrefAndHyperrefIfExist{definition:standard_order_on_natural_numbers}{standard order on $\bbN$}.

    The map $\psi: \bbN \to \bbZ_{> 0}$ defined by
    $$\psi(n) = [(n+1, 1)]$$
    is an isomorphism of totally ordered semirings. Consequently, $\bbZ_{> 0}$ is order-isomorphic to $\bbN$, and in particular $\bbZ_{> 0}$ is well-ordered.
\end{proposition}

\begin{proof}
\TODO{AI generated, verify}
    We verify each structure in turn.

    \begin{enumerate}
        \item $\psi$ is injective. If $\psi(n) = \psi(m)$ for $n,m \in \bbN$, then $[(n+1,1)] = [(m+1,1)]$ in $\bbZ$, so $(n+1,1) \sim (m+1,1)$ in $\bbN \times \bbN$. This means $n+1+1 = 1 + m+1$, i.e. $n+2 = m+2$. By \Cref{proposition:natural_numbers_addition_associative_commutative} (cancellativity), $n = m$.

        \item $\psi$ is surjective onto $\bbZ_{> 0}$. Let $x = [(a,b)] \in \bbZ_{> 0}$. Since $x > 0_\bbZ$, we have $a > b$ in $\bbN$. By the \CrefAndHyperrefIfExist{definition:standard_order_on_natural_numbers}{standard order on $\bbN$}, there exists $n \in \bbN$ such that $a = b + n$. Then:
        $$(a,b) = (b+n, b) \sim (n+1, 1),$$
        because $(b+n) + 1 = b + n + 1 = b + (n+1)$. Hence $x = [(n+1,1)] = \psi(n)$, so $\psi$ is surjective.

        \item $\psi$ preserves addition. For $n,m \in \bbN$,
        \begin{align*}
        \psi(n) + \psi(m) &= [(n+1, 1)] + [(m+1, 1)] = [(n+m+2,\, 2)], \\
        \psi(n+m) &= [(n+m+1,\, 1)].
        \end{align*}
        We check that $(n+m+2,2) \sim (n+m+1,1)$ in $\bbN \times \bbN$:
        $$(n+m+2) + 1 = n+m+3 = 2 + (n+m+1)$$
        by associativity and commutativity of addition in $\bbN$\CrefIfExists{proposition:natural_numbers_addition_associative_commutative}. Thus $\psi(n) + \psi(m) = \psi(n+m)$.

        \item $\psi$ preserves multiplication. For $n,m \in \bbN$,
        \begin{align*}
        \psi(n) \cdot \psi(m) &= [(n+1,1)] \cdot [(m+1,1)] = [((n+1)(m+1)+1\cdot 1,\; (n+1)\cdot 1 + 1\cdot(m+1))] \\
        &= [(nm+n+m+2,\; n+m+2)], \\
        \psi(nm) &= [(nm+1,\, 1)].
        \end{align*}
        We verify $(nm+n+m+2,\, n+m+2) \sim (nm+1,\, 1)$:
        $$(nm+n+m+2) + 1 = nm+n+m+3 = (n+m+2) + (nm+1)$$
        by associativity and commutativity of addition in $\bbN$\CrefIfExists{proposition:natural_numbers_addition_associative_commutative}. Hence $\psi(n) \cdot \psi(m) = \psi(nm)$.

        \item $\psi$ preserves unity. The multiplicative identity of $\bbN$ is $1$, and $\psi(1) = [(2,1)] = 1_\bbZ$\CrefIfExists{definition:integers_from_natural_numbers}.

        \item $\psi$ preserves order. For $n,m \in \bbN$, we have $n < m$ iff there exists $k \in \bbN$ with $n + k = m$ (\Cref{definition:standard_order_on_natural_numbers}). Then $\psi(n) = [(n+1,1)]$ and $\psi(m) = [(m+1,1)] = [(n+k+1,1)]$. In $\bbZ$, we have $\psi(m) - \psi(n) = [(n+k+1,1)] - [(n+1,1)] = [(k, 1)]$, which is positive because $k \geq 1$. Hence $\psi(n) < \psi(m)$. Conversely, if $\psi(n) < \psi(m)$ in $\bbZ$, then by the identification $\bbZ_{> 0} = \{[(a,b)] : a > b\}$ we have $n < m$. Thus $\psi$ is an order isomorphism.
    \end{enumerate}
\end{proof}
\input{../_concepts/theorem_integers_universal_property_ring_of_integers_is_the_initial_object_in_the_category_of_rings.tex}
\begin{theorem} \label{theorem:integers_uniqueness}
\TODO{AI generated, verify}
    Let $R$ be an \CrefAndHyperrefIfExist{definition:ordered_ring}{ordered ring} with positive cone $P$. Suppose that $P$ is well-ordered under the induced order. Let $\bbZ$ be the \CrefAndHyperrefIfExist{definition:ring_of_integers_from_natural_numbers}{ring of integers}. Then there exists a unique injective ring homomorphism $\phi: \bbZ \to R$ such that $\phi$ maps the positive cone $\bbZ_{> 0}$ into $P$.
\end{theorem}

\begin{proof}
\TODO{AI generated, verify}
    We construct $\phi$ using the \CrefAndHyperrefIfExist{theorem:integers_universal_property_ring_of_integers_is_the_initial_object_in_the_category_of_rings}{universal property of $\bbZ$}. Let $1_R \in R$ be the multiplicative identity of $R$. Since $R$ is an ordered ring, we have $1_R \in P$\CrefIfExists{definition:ordered_ring}. The universal property gives a unique ring homomorphism $\phi: \bbZ \to R$ with $\phi(1_\bbZ) = 1_R$. It remains to show that $\phi$ is injective and maps $\bbZ_{> 0}$ into $P$.

    \begin{enumerate}
        \item $\phi$ maps $\bbZ_{> 0}$ into $P$. By the \CrefAndHyperrefIfExist{proposition:positive_integers_correspond_to_natural_numbers}{positive integers correspondence}, every $x \in \bbZ_{> 0}$ can be written as $n \cdot 1_\bbZ$ for some $n \in \bbN$ (where the product is the $n$-fold sum). Since $\phi$ is a ring homomorphism, $\phi(n \cdot 1_\bbZ) = n \cdot 1_R$. We prove by induction on $n$ that $n \cdot 1_R \in P$. The base case $n = 1$ holds because $1 \cdot 1_R = 1_R \in P$ in an ordered ring. Suppose $k \cdot 1_R \in P$ for some $k \in \bbN$. Then $(k+1) \cdot 1_R = k \cdot 1_R + 1_R \in P$ because $P$ is closed under addition in an ordered ring\CrefIfExists{definition:ordered_ring}. Hence $\phi(\bbZ_{> 0}) \subseteq P$.

        \item $\phi$ is injective. Suppose $\phi(x) = 0_R$ for some $x \in \bbZ$. We show $x = 0_\bbZ$.

        If $x \in \bbZ_{> 0}$, then $x = n \cdot 1_\bbZ$ for some $n \in \bbN$, so $n \cdot 1_R = 0_R$. But in an ordered ring, $1_R > 0_R$, and by induction $n \cdot 1_R > 0_R$ for every $n \in \bbN$: the base case $n = 1$ holds by definition, and if $k \cdot 1_R > 0_R$, then $(k+1) \cdot 1_R = k \cdot 1_R + 1_R > 0_R$ because the sum of two positive elements is positive. Thus $n \cdot 1_R \neq 0_R$ for any $n \in \bbN$, a contradiction. Hence $x \notin \bbZ_{> 0}$.

        If $x \in \bbZ_{< 0}$, then $-x \in \bbZ_{> 0}$, so $\phi(-x) = -\phi(x) = 0_R$. By the same argument, $\phi(-x) = 0_R$ forces $x = 0_\bbZ$, a contradiction. Hence $x \notin \bbZ_{< 0}$.

        Therefore $x = 0_\bbZ$, so $\ker \phi = \{0_\bbZ\}$. Thus $\phi$ is injective.
    \end{enumerate}
\end{proof}
\begin{corollary} \label{corollary:integers_existence_uniqueness}
\TODO{AI generated, verify}
    The ring of integers $\bbZ$ exists and is unique up to isomorphism among ordered rings whose positive cone is well-ordered. Specifically:
    \begin{enumerate}
        \item (Existence) The construction in \CrefAndHyperrefIfExist{definition:integers_from_natural_numbers}{integers from natural numbers} yields an ordered ring satisfying the properties of \CrefAndHyperrefIfExist{definition:ring_of_integers}{ring of integers}. Equivalently, $\bbZ$ is (up to isomorphism) the \CrefAndHyperrefIfExist{definition:grothendieck_ring_of_a_commutative_semiring}{Grothendieck ring} of \CrefAndHyperrefIfExist{definition:natural_numbers_with_zero}{$\bbN_0$}\CrefIfExists{proposition:grothendieck_ring_of_natural_numbers_with_zero}.
        \item (Uniqueness) If $R$ is another ordered ring with well-ordered positive cone, then $R \cong \bbZ$ via the unique injective homomorphism of \CrefAndHyperrefIfExist{theorem:integers_uniqueness}{integers uniqueness theorem}. Moreover, the homomorphism is also surjective, hence an isomorphism.
    \end{enumerate}
\end{corollary}

\begin{proof}
\TODO{AI generated, verify}
    Existence. This follows from \Cref{proposition:integers_form_ring} and \Cref{proposition:integers_form_totally_ordered_ring}, which establish that the constructed $\bbZ$ is a totally ordered integral domain whose positive cone $\bbZ_{> 0} = \{[(a,b)] : a > b\}$ is well-ordered. The equivalence with the Grothendieck ring follows from \Cref{proposition:grothendieck_ring_of_natural_numbers_with_zero}.

    Uniqueness. Let $R$ be an ordered ring with positive cone $P$, where $P$ is well-ordered. By \Cref{theorem:integers_uniqueness}, there exists a unique injective homomorphism $\phi: \bbZ \to R$ mapping $\bbZ_{> 0}$ into $P$. We show that $\phi$ is surjective.

    We first prove that $1_R$ is the least element of $P$. Since $P$ is well-ordered, it has a least element; denote it by $\mu$. Because $1_R \in P$, we have $\mu \leq 1_R$. We claim $\mu = 1_R$.

    Since $\mu^2 \in P$ (because $P \cdot P \subseteq P$ in an ordered ring\CrefIfExists{definition:ordered_ring}), minimality of $\mu$ gives $\mu \leq \mu^2$, i.e.\ $\mu^2 - \mu \geq 0_R$. If $\mu < 1_R$, then $1_R - \mu \in P$ by trichotomy\CrefIfExists{definition:ordered_ring}. The product
    $$\mu(1_R - \mu) = \mu - \mu^2$$
    lies in $P$ because both factors are in $P$. But $\mu \leq \mu^2$ means $\mu - \mu^2 \leq 0_R$, so the only possibility is $\mu - \mu^2 = 0_R$, i.e.\ $\mu(1_R - \mu) = 0_R$.

    We claim that $R$ has no zero divisors. If $ab = 0_R$ with $a, b \neq 0_R$, then by trichotomy in the ordered ring\CrefIfExists{definition:ordered_ring} we may assume without loss of generality that $a, b \in P$ (replacing each element by its negation as needed). But then $ab \in P$ by compatibility of multiplication with the order, contradicting $0_R \notin P$. Hence $\mu(1_R - \mu) = 0_R$ with $\mu \in P$ and $\mu \neq 0_R$ implies $1_R - \mu = 0_R$, i.e.\ $\mu = 1_R$. This contradicts the assumption $\mu < 1_R$. Therefore $\mu = 1_R$.

    We now show that every element of $P$ is of the form $n \cdot 1_R$ for some $n \in \bbN$. Let $x \in P$. Consider the set
    $$T_x = \{x - n \cdot 1_R : n \in \bbN,\; x - n \cdot 1_R \geq 0_R\}.$$
    If $T_x \subseteq \{0_R\}$, then $x = n \cdot 1_R$ for some $n$, and we are done. Otherwise, $T_x \cap P$ is a nonempty subset of $P$, so it has a least element by well-ordering; denote this least element by $\rho$. We claim $\rho = 0_R$. If $\rho > 0_R$, then $\rho \in P$ and $\rho < 1_R$: for if $\rho \geq 1_R$, then $x - n \cdot 1_R \geq 1_R$ for some $n$, so $x - (n+1) \cdot 1_R = \rho - 1_R \geq 0_R$, meaning $\rho - 1_R \in T_x$. Since $\rho - 1_R < \rho$ and $\rho - 1_R \geq 0_R$, if $\rho - 1_R > 0_R$ it contradicts the minimality of $\rho$, and if $\rho - 1_R = 0_R$ then $\rho = 1_R$. But we just showed that $1_R$ is the least element of $P$, so no such $\rho \in P$ with $\rho < 1_R$ exists. Hence $\rho = 0_R$, meaning $x = n \cdot 1_R$ for some $n \in \bbN$.

    We finally establish the surjectivity of $\phi$. For any $r \in R$, by trichotomy in the ordered ring $R$\CrefIfExists{definition:ordered_ring}, exactly one of $r > 0_R$, $r = 0_R$, or $r < 0_R$ holds. If $r = 0_R$, then $r = \phi(0_\bbZ)$. If $r > 0_R$, then $r \in P$, so $r = n \cdot 1_R = \phi(n \cdot 1_\bbZ)$ for some $n \in \bbN$. If $r < 0_R$, then $-r \in P$, so $-r = n \cdot 1_R$ for some $n \in \bbN$, meaning $r = -(n \cdot 1_R) = \phi(-n \cdot 1_\bbZ)$. Hence $\phi$ is surjective.

    Since $\phi$ is bijective and a ring homomorphism, it is an isomorphism.
\end{proof}

\subsection{Ordered Ring Structure on \texorpdfstring{$\mathbb Z$}{Z}}

\begin{definition} \label{definition:standard_order_on_ring_of_integers}
Let $\bbZ$ be the \CrefAndHyperrefIfExist{definition:ring_of_integers}{ring of integers}, constructed as equivalence classes $[(a,b)]$ of pairs $(a,b) \in \bbN \times \bbN$ under the relation $(a,b) \sim (c,d) \iff a + d = b + c$.

Define a binary relation \hl{$\le$} on $\bbZ$ by
$$\hlin{[(a,b)] \le [(c,d)] \iff b + c \geq a + d,}$$
where the inequality $b + c \geq a + d$ is with respect to the \CrefAndHyperrefIfExist{definition:standard_order_on_natural_numbers}{standard order on $\bbN$}.

This relation is well-defined: if $(a,b) \sim (a',b')$ and $(c,d) \sim (c',d')$, then $a + b' = a' + b$ and $c + d' = c' + d$. Adding these equations gives
$$(b + c) + (b' + d') = (a + d) + (a' + d'),$$
so $b + c \geq a + d \iff b' + c' \geq a' + d'$ by cancellativity of addition on $\bbN$. Hence the truth value of $[(a,b)] \le [(c,d)]$ is independent of the chosen representatives.

The \hldef{standard order on $\bbZ$} is the pair $(\bbZ, \le)$. We write \hl{$x < y$} if $x \le y$ and $x \neq y$, and \hl{$y > x$} (resp. \hl{$y \ge x$}) if $x < y$ (resp. $x \le y$). The ring of integers becomes a \CrefAndHyperrefIfExist{definition:ordered_ring}{totally ordered ring} under this standard order\CrefIfExists{proposition:integers_form_totally_ordered_ring}.
\end{definition}

\begin{proposition} \label{proposition:integers_form_totally_ordered_ring}
Let $\bbZ$ be the \CrefAndHyperrefIfExist{definition:integers_from_natural_numbers}{ring of integers}\CrefIfExists{proposition:integers_form_ring}. Equipped with the \CrefAndHyperrefIfExist{definition:standard_order_on_ring_of_integers}{standard order on $\bbZ$}, $(\bbZ, +, \cdot, \le)$ is a \CrefAndHyperrefIfExist{definition:ordered_ring}{totally ordered ring}. Moreover, $\bbZ_{>0}$ is a \CrefAndHyperrefIfExist{definition:well_ordered_set}{well ordered set}, i.e. every nonempty subset of the positive cone $\bbZ_{> 0}$ has a least element.
\CrefIfExists{proposition:positive_integers_correspond_to_natural_numbers}
\end{proposition}

\begin{proof}
We verify each condition in the \CrefIfExists{definition:ordered_ring} step by step. Recall that for $x = [(a,b)]$ and $y = [(c,d)]$, we have $x \leq y \iff b + c \geq a + d$.

We first show that $\leq$ is a partial order on $\bbZ$. Let $x = [(a,b)], y = [(c,d)], z = [(e,f)] \in \bbZ$.
\begin{enumerate}
    \item The relation $\leq$ is reflexive because for any $x \in \bbZ$, we have $b + a = a + b$ by \Cref{proposition:natural_numbers_addition_associative_commutative}, so $b + a \geq a + b$ in $\bbN$, hence $x \leq x$.

    \item The relation $\leq$ is antisymmetric because if $x \leq y$ and $y \leq x$, then $b + c \geq a + d$ and $a + d \geq b + c$ in $\bbN$, so $b + c = a + d$. Thus $(a,b) \sim (c,d)$, meaning $x = y$.

    \item The relation $\leq$ is transitive because if $x \leq y$ and $y \leq z$, then $b + c \geq a + d$ and $d + e \geq c + f$ in $\bbN$. Adding these inequalities gives $(b + c) + (d + e) \geq (a + d) + (c + f)$, which rearranges to $(b + e) + (c + d) \geq (a + f) + (c + d)$ by \Cref{proposition:natural_numbers_addition_associative_commutative}. By cancellativity of addition on $\bbN$ (\Cref{proposition:natural_numbers_addition_associative_commutative}), $b + e \geq a + f$, hence $x \leq z$.
\end{enumerate}

Next we check translation invariance, i.e. that if $x \leq y$, then $x + z \leq y + z$ for all $z \in \bbZ$. Let $x = [(a,b)], y = [(c,d)], z = [(e,f)] \in \bbZ$. If $x \leq y$, then $b + c \geq a + d$. Now $x + z = [(a+e, b+f)]$ and $y + z = [(c+e, d+f)]$. We check:
$$(b+f) + (c+e) = (b+c) + (e+f) \geq (a+d) + (e+f) = (a+e) + (d+f),$$
where the inequality follows from $b + c \geq a + d$ and \Cref{proposition:natural_numbers_order_compatibility} (compatibility of order with addition on $\bbN$). Hence $x + z \leq y + z$.

We now establish compatibility with multiplication, i.e. that if $x \geq 0_\bbZ$ and $y \geq 0_\bbZ$, then $xy \geq 0_\bbZ$. First, we identify the non-negative elements. For $x = [(a,b)] \in \bbZ$, we have $x \geq 0_\bbZ = [(1,1)]$ iff $1 + a \geq b + 1$, i.e., $a \geq b$ in $\bbN$. Thus
$$\bbZ_{\geq 0} = \{[(a,b)] : a \geq b\}.$$

Now let $x, y \in \bbZ_{\geq 0}$ with $x = [(a,b)]$ and $y = [(c,d)]$, so $a \geq b$ and $c \geq d$. We must show $xy = [(ac+bd, ad+bc)] \geq 0_\bbZ$, i.e., $(ac+bd) + 1 \geq (ad+bc) + 1$, equivalently $ac + bd \geq ad + bc$.

If $a = b$ or $c = d$, then $x = 0_\bbZ$ or $y = 0_\bbZ$, so $xy = 0_\bbZ \geq 0_\bbZ$\CrefIfExists{lemma:multiplication_by_zero_in_a_ring_is_zero}.
Assume $a > b$ and $c > d$. By \Cref{definition:standard_order_on_natural_numbers}, there exist $r, s \in \bbN$ such that $a = b + r$ and $c = d + s$.
Expanding by distributivity on $\bbN$ (\Cref{proposition:natural_numbers_multiplication_associative_commutative_distributive}):
\begin{align*}
ac + bd &= (b+r)(d+s) + bd = bd + bs + rd + rs + bd, \\
ad + bc &= b(d+s) + (b+r)d = bd + bs + bd + rd.
\end{align*}
Thus $ac + bd = ad + bc + rs$. Since $r, s \in \bbN$, we have $rs \in \bbN$, so $ac + bd \geq ad + bc$ in $\bbN$ by \Cref{definition:standard_order_on_natural_numbers}. Hence $xy \geq 0_\bbZ$.

We verify trichotomy, i.e. that for every $x \in \bbZ$, exactly one of $x > 0_\bbZ$, $x = 0_\bbZ$, or $x < 0_\bbZ$ holds. For any $x = [(a,b)] \in \bbZ$, by \Cref{proposition:natural_numbers_order_compatibility} (trichotomy of order on $\bbN$), exactly one of $a > b$, $a = b$, or $a < b$ holds.
\begin{enumerate}
    \item If $a > b$, then $1 + a > 1 + b$, so $x > 0_\bbZ$.
    \item If $a = b$, then $(a,b) \sim (1,1)$ because $a + 1 = b + 1$, so $x = 0_\bbZ$.
    \item If $a < b$, then $[-x] = [(b,a)] > 0_\bbZ$ by case (i), meaning $x < 0_\bbZ$.
\end{enumerate}
Thus, we have now proven that $\bbZ$ is a totally ordered ring under the standard order.

Finally, we show that $\bbZ_{> 0}$ is well-ordered. Let $S \subseteq \bbZ_{> 0}$ be nonempty. For each $x = [(a,b)] \in S$, we have $x > 0_\bbZ$, so $a > b$ in $\bbN$. Define the difference of $x$ as $\operatorname{diff}(x) = r$, where $r \in \bbN$ is the unique element such that $a = b + r$. This is well-defined because if $(a,b) \sim (a',b')$ with both representing positive elements, then $a + b' = a' + b$, and since $a = b + r$ and $a' = b' + r'$, we get $(b+r) + b' = (b'+r') + b$, so $r = r'$ by \Cref{proposition:natural_numbers_addition_associative_commutative} (cancellativity).

Let $D = \{\operatorname{diff}(x) : x \in S\} \subseteq \bbN$. Since $S$ is nonempty, so is $D$. By \Cref{proposition:natural_numbers_well_ordered}, $D$ has a least element $d_0$. Pick any $x_0 \in S$ with $\operatorname{diff}(x_0) = d_0$.

For any $y \in S$, let $d_y = \operatorname{diff}(y) \in D$. Since $d_0 \leq d_y$ in $\bbN$, we have two cases:
\begin{enumerate}
    \item If $d_0 = d_y$, write $x_0 = [(a_0, b_0)]$ and $y = [(c, d)]$ with $a_0 = b_0 + d_0$ and $c = d + d_0$. Then $(a_0, b_0) \sim (c, d)$ because $a_0 + d = (b_0 + d_0) + d = b_0 + (d + d_0) = b_0 + c$ by \Cref{proposition:natural_numbers_addition_associative_commutative}. Hence $x_0 = y$.
    \item If $d_0 < d_y$, then $d_y = d_0 + s$ for some $s \in \bbN$. Write $x_0 = [(b_0 + d_0, b_0)]$ and $y = [(d + d_0 + s, d)]$. Then
    $$y - x_0 = [(d + d_0 + s + b_0, d + b_0 + d_0)].$$
    The first component exceeds the second by $s \in \bbN$, so $y - x_0 > 0_\bbZ$, meaning $x_0 < y$.
\end{enumerate}

Thus $x_0 \leq y$ for all $y \in S$, making $x_0$ the least element.
\end{proof}

\begin{definition} \label{definition:ring_of_integers}
    The \hldef{ring of integers}, universally denoted by \hl{$\bbZ$}, is the unique (up to \CrefAndHyperref{definition:ring_homomorphism}{isomorphism}) \CrefAndHyperref{definition:ordered_ring}{ordered}\CrefIfExists{proposition:integers_form_ordered_ring} \CrefAndHyperrefIfExist{definition:domain_integral_domain}{integral domain} whose \CrefAndHyperrefIfExist{definition:positive_cone_in_an_ordered_monoid}{positive cone} $\bbZ_{>0}$ is non-empty and \CrefAndHyperrefIfExist{definition:well_ordered_set}{well-ordered}, i.e. every non-empty subset of $\bbZ_{>0}$ has a least element\CrefIfExists{corollary:integers_existence_uniqueness}.
    % such that there exists a subset $P \subset \bbZ$ such that 
    % \begin{enumerate}
    %     \item $P$ is nonempty, and
    %     \item Every non-empty subset of $P$ contains a least element.
    % \end{enumerate}
    \TextIfExists{definition:category_of_rings}{
        The ring $\bbZ$ is also \CrefAndHyperrefIfExist{theorem:integers_universal_property_ring_of_integers_is_the_initial_object_in_the_category_of_rings}{uniquely characterized} as the \CrefAndHyperrefIfExist{definition:initial_final_zero_objects_of_a_category}{initial object} in the \CrefAndHyperrefIfExist{definition:category_of_rings}{category of rings}
    }
    \TextIfExists{definition:ring_of_integers_from_natural_numbers}{
        See also \Cref{definition:ring_of_integers_from_natural_numbers} for a construction of the ring of integers.
    }
\end{definition}
\begin{definition}[Partially ordered set] \label{definition:partially_ordered_set}
    \begin{enumerate}
        \item 
        A \hldef{partially ordered set} (or \hldef{poset}), or \hldef{ordered set} is a pair $(P, \leq)$ where $P$ is a set and 
        \[
        \leq : P \times P \to \{\text{true}, \text{false}\}
        \]
        is a \CrefAndHyperrefIfExist{definition:n_ary_relation_on_sets}{binary relation} on $P$ satisfying the following axioms for all $a,b,c \in P$:
        \begin{itemize}
            \item \CrefAndHyperrefIfExist{definition:reflexive_symmetric_antisymmetric_transitive_total_binary_relations_on_a_set}{Reflexivity}: $a \leq a$,
            \item \CrefAndHyperrefIfExist{definition:reflexive_symmetric_antisymmetric_transitive_total_binary_relations_on_a_set}{Antisymmetry}: if $a \leq b$ and $b \leq a$, then $a = b$,
            \item \CrefAndHyperrefIfExist{definition:reflexive_symmetric_antisymmetric_transitive_total_binary_relations_on_a_set}{Transitivity}: if $a \leq b$ and $b \leq c$, then $a \leq c$.
        \end{itemize}
        The relation $\leq$ is called an \hldef{order} or a \hldef{partial order}

        One writes \hl{$a \geq b$} synonymously for $b \leq a$. One might also synonymously write \hl{$a < b$} and \hl{$a > b$} respectively for 
        \begin{itemize}
            \item $a \leq b$ and $a \neq b$,
            \item $a \geq b$ and $a \neq b$.
        \end{itemize}
        This binary relation $<$ is then a \CrefAndHyperrefIfExist{definition:strict_partially_ordered_set}{strict partial order} on $P$\CrefIfExists{proposition:correspondence_between_partial_orders_and_strict_partial_orders_on_a_set}.

        One also says ``\hldef{$a$ is less than or equal or $b$}'', ``\hldef{$a$ is at most $b$}'', or ``\hldef{$a$ is bounded above by $b$}'' to mean $a \leq b$, says ``\hldef{$a$ is greater than or equal to $b$}'', ``\hldef{$a$ is at least $b$}'', or ``\hldef{$a$ is bounded below by $b$}'' to mean $a \geq b$, says ``\hldef{$a$ is less than $b$}'' or ``\hldef{$a$ is strictly bounded above by $b$}'' to mean $a < b$, and says ``\hldef{$a$ is greater than $b$}'' or ``\hldef{$a$ is strictly bounded below by $b$}'' to mean $a > b$.

        \item A partially ordered set $(P, \leq)$ is called a \hldef{directed partially ordered set} if for every pair $a,b \in P$, there exists $c \in P$ such that
        \[
        a \leq c \quad \text{and} \quad b \leq c.
        \]

        \item A partially ordered set $(P, \leq)$ is called a \hldef{codirected partially ordered set} (or \hldef{downward directed poset}) if for every pair $a,b \in P$, there exists $d \in P$ such that
        \[
        d \leq a \quad \text{and} \quad d \leq b.
        \]
        \end{enumerate}
\end{definition}

\begin{definition}[Total order on a set] \label{definition:total_order_linear_order_on_a_set}
A \hldef{total order} (or \hldef{linear order}) on a set $S$ is a \CrefAndHyperrefIfExist{definition:n_ary_relation_on_sets}{binary relation} $\leq$ such that $(S, \leq)$ is a \CrefAndHyperrefIfExist{definition:partially_ordered_set}{partially ordered set} in which any two elements are comparable. That is, for all $a, b \in S$, either $a \leq b$ or $b \leq a$.

A set equipped with a total order is called a \hldef{totally ordered set}.
\TextIfExists{definition:partially_ordered_set}{
    In particular, all totally ordered sets are partially ordered sets.
}
\end{definition}
\begin{definition} \label{definition:ordered_monoid}
An \hldef{ordered monoid} is a \CrefAndHyperrefIfExist{definition:monoid}{monoid} $(M, \cdot, e)$ equipped with a \CrefAndHyperrefIfExist{definition:partially_ordered_set}{partial order} $\le$ that is compatible with the monoid operation. That is, for all $a, b, c \in M$, if $a \le b$, then:
$$a \cdot c \le b \cdot c \quad \text{and} \quad c \cdot a \le c \cdot b.$$

If the order $\le$ is a \CrefAndHyperrefIfExist{definition:total_order_linear_order_on_a_set}{total order}, then $M$ is called a \hldef{totally ordered monoid}.

If the monoid structure is commutative/abelian, then we call $M$ is called a \hldef{commutative/abelian ordered monoid}.
\end{definition}
\begin{definition}[Ordered abelian group] \label{definition:ordered_abelian_group}
    \TODO{AI generated, verify}
    \begin{enumerate}
        \item An \hldef{ordered abelian group} is an \CrefAndHyperrefIfExist{definition:group}{abelian group} $(G,+)$ equipped with a \CrefAndHyperrefIfExist{definition:partially_ordered_set}{partial order} $\leq$ that is compatible with the group operation. That is, for all $a,b,c \in G$, if $a \leq b$, then
        $$a + c \leq b + c.$$
        If the partial order $\leq$ is a \CrefAndHyperrefIfExist{definition:total_order_linear_order_on_a_set}{total order}, then $G$ is called a \hldef{totally ordered abelian group}.

        \item Equivalently, an ordered abelian group may be defined as an abelian group $(G,+)$ equipped with a \CrefAndHyperrefIfExist{definition:strict_partially_ordered_set}{strict partial order} $<$ that is compatible with the group operation: for all $a,b,c \in G$, if $a < b$, then
        $$a + c < b + c.$$
        These two definitions are identified via the natural correspondence between a partial order $\leq$ and its associated strict partial order $<$ given by $a < b \iff a \leq b \text{ and } a \neq b$\CrefIfExists{proposition:correspondence_between_partial_orders_and_strict_partial_orders_on_a_set}. When using the strict formulation, we adopt the notations that $a \leq b$ means $a < b$ or $a = b$, and \hl{$b > a$} (resp. \hl{$b \geq a$}) means $a < b$ (resp. $a \leq b$).
    \end{enumerate}

    Writing $0 \in G$ for the identity of $G$, elements $a \in G$ such that $a > 0$ are called \hldef{positive}, elements $a \in G$ such that $a < 0$ are called \hldef{negative}, elements $a \in G$ such that $a \geq 0$ are called \hldef{non-negative}, and elements $a \in G$ such that $a \leq 0$ are called \hldef{non-positive}.
\end{definition}
\begin{definition} \label{definition:subtraction_of_an_element_of_a_set_by_a_uniquely_invertible_element_for_a_commutative_binary_operation_written_additively}
    \TODO{AI generated, verify}
    Let $+$ be a \CrefAndHyperref{definition:commutative_binary_operation_on_a_set}{commutative} \CrefAndHyperref{definition:n_ary_operation_on_a_set}{binary operation} on a \CrefAndHyperref{definition:zermelo_fraenkel_set_theory}{set} $S$, written additively, with an \CrefAndHyperref{definition:left_right_identity_element_of_a_binary_operation_on_a_set}{identity element} \hl{$0 \in S$}. Let $x,y \in S$ be elements such that $y$ has a unique \CrefAndHyperref{definition:left_right_inverse_of_an_element_for_a_binary_operation_on_a_set}{inverse} (guaranteed when $(S,+)$ is a commutative group), denoted by \hl{$-y$}.

    The \hldef{subtraction of $y$ from $x$} is defined as
    $$\hlin{x - y := x + (-y)}.$$
    The element $x - y$ is called the \hldef{difference of $x$ and $y$}.
\end{definition}

\begin{definition}[Ordered ring] \label{definition:ordered_ring}
    \begin{enumerate}
        \item An \hldef{ordered ring} is a \CrefAndHyperrefIfExist{definition:ring}{ring} $(R,+,\cdot)$ equipped with a \CrefAndHyperrefIfExist{definition:partially_ordered_set}{partial order} $\leq$ such that
        \begin{enumerate}
            \item the additive group $(R,+)$ is an \CrefAndHyperrefIfExist{definition:ordered_abelian_group}{ordered abelian group}, i.e.\ for all $a,b,c \in R$, if $a \leq b$, then $a + c \leq b + c$;
            \item the order is compatible with multiplication: if $a \geq 0$ and $b \geq 0$, then $ab \geq 0$.
        \end{enumerate}
        The set of \hldef{non-negative elements} of $R$ is denoted by \hl{$R_{\ge 0} = \{a \in R : a \ge 0\}$}.

        If the partial order $\leq$ is a \CrefAndHyperrefIfExist{definition:total_order_linear_order_on_a_set}{total order}, then $R$ is called a \hldef{totally ordered ring}. 
        Equivalently, an ordered ring is totally ordered if and only if it satisfies \hldef{trichotomy}: for every $a \in R$, exactly one of $a > 0$, $a = 0$, or $a < 0$ holds.

        \item Equivalently, an ordered ring may be defined as a ring $(R,+,\cdot)$ equipped with a \CrefAndHyperrefIfExist{definition:strict_partially_ordered_set}{strict partial order} $<$ such that
        \begin{enumerate}
            \item the additive group $(R,+)$ is an ordered abelian group under $<$ (i.e.\ for all $a,b,c \in R$, if $a < b$, then $a + c < b + c$);
            \item the order is compatible with multiplication in the following sense: for all $a,b \in R$, if $a > 0$ and $b > 0$, then $ab > 0$.
        \end{enumerate}
    \end{enumerate}

    These two definitions are identified via the natural correspondence between a partial order $\leq$ and its associated strict partial order $<$ given by $a < b \iff a \leq b \text{ and } a \neq b$\CrefIfExists{proposition:correspondence_between_partial_orders_and_strict_partial_orders_on_a_set}. When using the strict formulation, we adopt the notations that $a \leq b$ means $a < b$ or $a = b$, and \hl{$b > a$} (resp.\ \hl{$b \geq a$}) means $a < b$ (resp.\ $a \leq b$). The set of \hldef{positive elements} is denoted by \hl{$R_{> 0} = \{a \in R : a > 0\}$}. For any $a,b \in R$, we have
    $$a - b \in R_{> 0} \iff a > b.$$
    \CrefIfExists{definition:subtraction_of_an_element_of_a_set_by_a_uniquely_invertible_element_for_a_commutative_binary_operation_written_additively}
\end{definition}
\begin{definition} \label{definition:standard_absolute_value_of_an_ordered_ring}
    Let $R$ be an \CrefAndHyperrefIfExist{definition:ordered_ring}{ordered ring} with \CrefAndHyperrefIfExist{definition:positive_cone_in_an_ordered_monoid}{positive cone} $R_{\geq 0}$.
    The \hldef{(standard) absolute value function} is the function
    $$|\cdot|: R \to R_{\geq 0}$$
    defined as follows:
    $$|r| = \begin{cases} r &\text{if } r \in P \\ -r &\text{if } -r \in P \\ 0 &\text{otherwise}.\end{cases}$$
    If $R$ is totally ordered, then it is indeed an \CrefAndHyperrefIfExist{definition:absolute_value_on_a_ring_valued_in_an_ordered_semiring}{absolute value} on $R$ valued in $R_{\geq 0}$. \CrefIfExists{lemma:absolute_value_function_of_an_ordered_ring_is_indeed_an_absolute_value} For a merely partially ordered ring, $|\cdot|$ may fail positive definiteness (if $r$ is incomparable to $0_R$, then $|r| = 0_R$ while $r \neq 0_R$) and multiplicativity.
\end{definition}
\begin{lemma}[Standard absolute value on a totally ordered ring is an absolute value] \label{lemma:absolute_value_function_of_an_ordered_ring_is_indeed_an_absolute_value}
    Let $R$ be a \CrefAndHyperrefIfExist{definition:ordered_ring}{totally ordered commutative ring} with \CrefAndHyperrefIfExist{definition:positive_cone_in_an_ordered_monoid}{positive cone} $R_{\geq 0}$.
    The \CrefAndHyperrefIfExist{definition:standard_absolute_value_of_an_ordered_ring}{absolute value function}
    $$|\cdot|: R \to R_{\geq 0}$$
    is an \CrefAndHyperrefIfExist{definition:absolute_value_on_a_ring_valued_in_an_ordered_semiring}{absolute value} on $R$ valued in the \CrefAndHyperrefIfExist{definition:ordered_semiring}{ordered semiring} $R_{\geq 0}$\CrefIfExists{proposition:positive_cone_of_ordered_ring_is_ordered_semiring}.
\end{lemma}

\begin{proof}
    Since $R$ is totally ordered, trichotomy holds: for every element $a \in R$, exactly one of $a > 0_R$, $a = 0_R$, or $-a > 0_R$ holds\CrefIfExists{definition:ordered_ring}. Hence the absolute value simplifies to
    $$|a| = \begin{cases} a &\text{if } a \geq 0_R, \\ -a &\text{if } a < 0_R.\end{cases}$$

    We verify the three axioms of an absolute value\CrefIfExists{definition:absolute_value_on_a_ring_valued_in_an_ordered_semiring}.

    \begin{enumerate}
        \item Positive definiteness: If $a = 0_R$, then $|a| = |0_R| = 0_R$ by the first case. Conversely, if $|a| = 0_R$, then either $a = 0_R$ (first case) or $-a = 0_R$ (second case). In the second case, $-a = 0_R$ implies $a = -0_R = 0_R$. Hence $|a| = 0_R$ if and only if $a = 0_R$.

        \item Multiplicativity: We must show $|ab| = |a||b|$ for all $a,b \in R$. We consider cases based on the sign of $a$ and $b$.
        \begin{enumerate}
            \item If $a \geq 0_R$ and $b \geq 0_R$, then $ab \geq 0_R$ by compatibility of the order with multiplication\CrefIfExists{definition:ordered_ring}. Hence $|ab| = ab = |a||b|$.
            \item If $a < 0_R$ and $b < 0_R$, then $-a > 0_R$ and $-b > 0_R$, so $(-a)(-b) > 0_R$ by compatibility. Since $(-a)(-b) = ab$ in any ring, we have $ab > 0_R$. Hence $|ab| = ab = (-a)(-b) = |a||b|$.
            \item If $a \geq 0_R$ and $b < 0_R$, then $-b > 0_R$. Since both $a$ and $-b$ are nonnegative, their product $a(-b) = -ab$ is nonnegative by compatibility, so $ab \leq 0_R$. Hence either $ab = 0_R$ (in which case $|ab| = 0_R = |a||b|$) or $ab < 0_R$ (in which case $|ab| = -ab = a(-b) = |a||b|$).
            \item If $a < 0_R$ and $b \geq 0_R$, the argument is symmetric: since both $-a$ and $b$ are nonnegative, their product $(-a)b = -(ab)$ is nonnegative, so $ab \leq 0_R$. Hence $|ab| = -ab = (-a)b = |a||b|$ or $|ab| = 0_R = |a||b|$ if $ab = 0_R$.
        \end{enumerate}

        \item Triangle inequality: We first show that for all $a \in R$, we have the two-sided bound $-|a| \leq a \leq |a|$. If $a \geq 0_R$, then $|a| = a$ and $-|a| = -a \leq 0_R \leq a = |a|$, so both inequalities hold. If $a < 0_R$, then $|a| = -a > 0_R$ and $a = -(-a) = -|a|$, so again $-|a| \leq a \leq |a|$.

        Now apply this bound to both $a$ and $b$ and add the resulting inequalities (valid by translation-invariance\CrefIfExists{definition:ordered_ring}):
        $$-(|a| + |b|) \leq a + b \leq |a| + |b|.$$

        Consider the sign of $a+b$. If $a+b \geq 0_R$, then $|a+b| = a+b \leq |a| + |b|$ by the right-hand inequality above. If $a+b < 0_R$, then $|a+b| = -(a+b)$. From the left-hand inequality, we have $-(|a| + |b|) \leq a + b$; negating both sides (which reverses the inequality in an ordered ring) gives $-(a+b) \leq |a| + |b|$, i.e., $|a+b| \leq |a| + |b|$. In both cases, $|a+b| \leq |a| + |b|$.
    \end{enumerate}

    This completes the verification that $|\cdot|$ satisfies all three axioms of an absolute value on $R$ valued in $R_{\geq 0}$.
\end{proof}
\begin{lemma}
    Let $R$ be an \CrefAndHyperrefIfExist{definition:ordered_ring}{ordered ring}.
    For $a,b,c,d \in R$, if $a < b$ and $c < d$, then $a + c < b+d$.
\end{lemma}
\begin{proof}
    \TODO{conitnue}
    Let $a,b,c,d \in R$ satisfy $a < b$ and $c < d$. Equivalently, $b-a \in P$ and $d-c \in P$. Since $P$ is closed under addition, we have
    $$(b-a) + (d-c) \in P.$$
    \CrefIfExists{definition:subtraction_of_an_element_of_a_set_by_a_uniquely_invertible_element_for_a_commutative_binary_operation_written_additively}
\end{proof}


% ============================================================
% Section 4: Divisibility
% ============================================================
\section{Divisibility}

\begin{definition} \label{definition:gcd_and_lcm_of_integers}
    \TODO{AI generated, verify}
    Let $a, b \in \bbZ$, not both zero.
    \begin{enumerate}
        \item A \hldef{greatest common divisor of $a$ and $b$} is an integer $d \geq 0$ such that
            \begin{enumerate}
                \item $d \mid a$ and $d \mid b$ (\CrefAndHyperrefIfExist{definition:divisibility_of_elements_of_rings}{divisibility});
                \item for any $c \in \bbZ$, if $c \mid a$ and $c \mid b$, then $c \mid d$.
            \end{enumerate}
        The greatest common divisor is unique and is denoted by \hl{$\gcd(a, b)$}.
        \item A \hldef{least common multiple of $a$ and $b$} is an integer $m \geq 0$ such that
            \begin{enumerate}
                \item $a \mid m$ and $b \mid m$;
                \item for any $n \in \bbZ$, if $a \mid n$ and $b \mid n$, then $m \mid n$.
            \end{enumerate}
        The least common multiple is unique and is denoted by \hl{$\operatorname{lcm}(a, b)$}.
    \end{enumerate}
\end{definition}

\begin{theorem}[Division Algorithm] \label{theorem:division_algorithm_for_integers_euclidean_algorithm_for_integers}
Let $a \in \bbZ$ and $b \in \bbZ_{> 0}$. Then there exist unique integers $q, r \in \bbZ$ such that
$$ a = bq + r \quad \text{and} \quad 0 \le r < b. $$
The integer $q$ is called the \hldef{quotient} and $r$ is called the \hldef{remainder} of the division of $a$ by $b$.
\end{theorem}
\begin{proposition} \label{proposition:bezout_identity_for_gcd_of_integers}
    \TODO{AI generated, verify}
    Let $a, b \in \bbZ$, not both zero. There exist integers $x, y$ such that
    $$\gcd(a, b) = ax + by.$$
\end{proposition}



% ============================================================
% Section 5: Congruences
% ============================================================
\section{Congruences}

\begin{definition} \label{definition:congruence_modulo_n}
    Let $n \in \bbZ_{>0}$. For $a, b \in \bbZ$, we say that $a$ and $b$ are \hldef{congruent modulo $n$} if $n \mid (a - b)$. This is denoted by
    $$\hlin{a \equiv b \pmod{n}.}$$
    If $n \nmid (a - b)$, we write $a \not\equiv b \pmod{n}$.
\end{definition}

\begin{definition} \label{definition:residue_class_modulo_n}
    \TODO{AI generated, verify}
    Let $n \in \bbZ_{>0}$ and let $a \in \bbZ$. The \hldef{residue class of $a$ modulo $n$} is the set
    $$\hlin{\overline{a} = \{ b \in \bbZ : b \equiv a \pmod{n} \}}$$
    where $\equiv \pmod{n}$ denotes \CrefAndHyperref{definition:congruence_modulo_n}{congruence modulo $n$}.

    The set of all residue classes modulo $n$ forms
    $$\hlin{\bbZ / n\bbZ = \{ \overline{0}, \overline{1}, \dots, \overline{n-1} \}}$$
    which is a finite \CrefAndHyperrefIfExist{definition:commutative_ring}{commutative ring} with $n$ elements under the operations
    $$\overline{a} + \overline{b} = \overline{a + b} \quad \text{and} \quad \overline{a} \cdot \overline{b} = \overline{ab}.$$
\end{definition}

\begin{lemma}[Number of Residue Classes Modulo m]
\label{lemma:number_of_residue_classes_modulo_m_is_m}
\TODO{AI generated, verify}
Let $m \in \bbZ_{> 0}$. The set of residue classes modulo $m$ is
$$ \bbZ / m\bbZ = \{ \overline{0}, \overline{1}, \dots, \overline{m-1} \}, $$
and these $m$ classes are distinct. Thus there are exactly $m$ residue classes modulo $m$.

Equivalently, the \CrefAndHyperrefIfExist{definition:residue_class_modulo_n}{residue classes} modulo $m$ form a partition of $\bbZ$ into exactly $m$ equivalence classes, each of the form $\overline{a} = \{ a + mk \mid k \in \bbZ \}$ for a unique $a \in \{0, 1, \dots, m-1\}$.
\end{lemma}
\begin{theorem} \label{theorem:chinese_remainder_theorem_for_pairwise_coprime_moduli}
    \TODO{AI generated, verify}
    Let $n_1, \dots, n_k \in \bbZ_{>0}$ be pairwise coprime integers, i.e., $\gcd(n_i, n_j) = 1$ for all $i \neq j$. Let $a_1, \dots, a_k \in \bbZ$. There exists an integer $x$ such that
    $$x \equiv a_i \pmod{n_i} \quad \text{for each } i = 1, \dots, k.$$
    Moreover, $x$ is unique modulo $n_1 \cdots n_k$. That is, if $y \in \bbZ$ also satisfies $y \equiv a_i \pmod{n_i}$ for all $i$, then
    $$x \equiv y \pmod{n_1 \cdots n_k}.$$
\end{theorem}



% ============================================================
% Section 6: Primes and Factorization
% ============================================================
\section{Primes and Factorization}

\begin{definition} \label{definition:prime_number}
    A \hldef{prime number} is a positive element of \CrefAndHyperrefIfExist{definition:ring_of_integers}{$\bbZ$} that is an \CrefAndHyperref{definition:irreducible_element_of_a_ring}{irreducible element} of $\bbZ$. 

    Explicitly, a prime number $p$ is a positive integer that is not $1$ and whose only positive integer \CrefAndHyperref{definition:divisibility_of_elements_of_rings}{divisors} are $1$ and $p$
    
    Equivalently, a prime number is a positive element of $\bbZ$ that is a \CrefAndHyperref{definition:prime_element_of_a_ring}{prime element} of $\bbZ$.
\end{definition}
\begin{theorem} \label{theorem:fundamental_theorem_of_arithmetic}
    \TODO{AI generated, verify}
    Every integer $n \geq 2$ can be written as a product of \CrefAndHyperref{definition:prime_number}{prime numbers}:
    $$n = p_1 p_2 \cdots p_k,$$
    where $p_1 \leq p_2 \leq \cdots \leq p_k$ are primes. This factorization is unique up to the order of the factors. That is, if also
    $$n = q_1 q_2 \cdots q_m$$
    with $q_1 \leq q_2 \leq \cdots \leq q_m$ primes, then $k = m$ and $p_i = q_i$ for all $i = 1, \dots, k$.
\end{theorem}



% ============================================================
% Section 7: Euler's Totient Function
% ============================================================
\section{Euler's Totient Function}

\begin{definition} \label{definition:euler_totient_function}
    \TODO{AI generated, verify}
    Let $n \in \bbZ_{>0}$. The \hldef{Euler totient function} evaluated at $n$, denoted by \hl{$\varphi(n)$}, is the number of integers $a$ with $1 \leq a \leq n$ such that $\gcd(a, n) = 1$:
    $$\hlin{\varphi(n) = \left| \{ a \in \bbZ : 1 \leq a \leq n,\; \gcd(a, n) = 1 \} \right|.}$$
    Equivalently, $\varphi(n)$ is the cardinality of the multiplicative group $(\bbZ / n\bbZ)^\times$.
\end{definition}

\begin{proposition} \label{proposition:formula_for_euler_totient_function}
    Let $n \in \bbZ_{>0}$ with \CrefAndHyperrefIfExist{theorem:fundamental_theorem_of_arithmetic}{prime factorization} $n = p_1^{e_1} \cdots p_k^{e_k}$, where $p_1, \dots, p_k$ are distinct \CrefAndHyperrefIfExist{definition:prime_number}{prime numbers} and $e_i \geq 1$. Then
    $$\varphi(n) = n \prod_{i=1}^k \left(1 - \frac{1}{p_i}\right).$$
    In particular, if $p$ is prime and $e \geq 1$, then $\varphi(p^e) = p^e - p^{e-1}$.
\end{proposition}



% ============================================================
% Section 8: Multiplicative Order and Primitive Roots
% ============================================================
\section{Multiplicative Order and Primitive Roots}

\begin{definition} \label{definition:multiplicative_order_of_integer_modulo_n}
    \TODO{AI generated, verify}
    Let $n \in \bbZ_{>0}$ and let $a \in \bbZ$ with $\gcd(a, n) = 1$. The \hldef{multiplicative order of $a$ modulo $n$}, denoted by \hl{$\operatorname{ord}_n(a)$}, is the smallest positive integer $d$ such that
    $$a^d \equiv 1 \pmod{n}.$$
\end{definition}

\begin{definition} \label{definition:primitive_root_modulo_n}
    \TODO{AI generated, verify}
    Let $n \in \bbZ_{>0}$ and let $a \in \bbZ$ with $\gcd(a, n) = 1$. We say that $a$ is a \hldef{primitive root modulo $n$} if the \CrefAndHyperref{definition:multiplicative_order_of_integer_modulo_n}{multiplicative order} of $a$ modulo $n$ equals $\varphi(n)$:
    $$\operatorname{ord}_n(a) = \varphi(n).$$
    Equivalently, $a$ is a primitive root modulo $n$ if the residue class $\overline{a}$ generates the multiplicative group $(\bbZ / n\bbZ)^\times$.
\end{definition}

\begin{theorem} \label{theorem:fermats_little_theorem}
    \TODO{AI generated, verify}
    Let $p$ be a \CrefAndHyperref{definition:prime_number}{prime number} and let $a \in \bbZ$ with $p \nmid a$. Then
    $$a^{p-1} \equiv 1 \pmod{p}.$$
    Equivalently, for any $a \in \bbZ$, $a^p \equiv a \pmod{p}$.
\end{theorem}

\begin{theorem} \label{theorem:eulers_theorem}
    \TODO{AI generated, verify}
    Let $n \in \bbZ_{>0}$ and let $a \in \bbZ$ with $\gcd(a, n) = 1$. Then
    $$a^{\varphi(n)} \equiv 1 \pmod{n}.$$
\end{theorem}



% ============================================================
% Section 9: Quadratic Residues
% ============================================================
\section{Quadratic Residues}

\begin{definition} \label{definition:legendre_symbol}
    Let $p$ be an odd \CrefAndHyperrefIfExist{definition:prime_number}{prime number} and let $a \in \bbZ$. The \hldef{Legendre symbol} is defined as
    $$\hlin{\left(\frac{a}{p}\right) = \begin{cases} 0 & \text{if } p \mid a, \\ 1 & \text{if } a \equiv x^2 \pmod{p} \text{ for some } x \in \bbZ \text{ with } p \nmid x, \\ -1 & \text{if } a \not\equiv x^2 \pmod{p} \text{ for all } x \in \bbZ.  \end{cases}}$$
    When $\left(\frac{a}{p}\right) = 1$, we say that $a$ is a \hldef{quadratic residue modulo $p$}. When $\left(\frac{a}{p}\right) = -1$, we say that $a$ is a \hldef{quadratic nonresidue modulo $p$}.
\end{definition}

\begin{definition} \label{definition:jacobi_symbol}
    \TODO{AI generated, verify}
    Let $n \in \bbZ_{>0}$ be odd with prime factorization $n = p_1^{e_1} \cdots p_k^{e_k}$, where $p_1, \dots, p_k$ are distinct odd \CrefAndHyperref{definition:prime_number}{prime numbers}. For any $a \in \bbZ$, the \hldef{Jacobi symbol} is defined as
    $$\hlin{\left(\frac{a}{n}\right) = \prod_{i=1}^k \left(\frac{a}{p_i}\right)^{e_i}},$$
    where each $\left(\frac{a}{p_i}\right)$ is the \CrefAndHyperref{definition:legendre_symbol}{Legendre symbol}.
\end{definition}

\begin{proposition} \label{proposition:properties_of_legendre_and_jacobi_symbols}
    \TODO{AI generated, verify}
    Let $p$ be an odd \CrefAndHyperref{definition:prime_number}{prime number} and let $a, b \in \bbZ$.
    \begin{enumerate}
        \item (Multiplicativity) $\left(\frac{ab}{p}\right) = \left(\frac{a}{p}\right)\left(\frac{b}{p}\right).$
        \item (Congruence) If $a \equiv b \pmod{p}$, then $\left(\frac{a}{p}\right) = \left(\frac{b}{p}\right).$
        \item (Supplementary law I) $\left(\frac{-1}{p}\right) = (-1)^{(p-1)/2}.$
        \item (Supplementary law II) $\left(\frac{2}{p}\right) = (-1)^{(p^2-1)/8}.$
    \end{enumerate}
\end{proposition}

\begin{theorem} \label{theorem:law_of_quadratic_reciprocity}
    \TODO{AI generated, verify}
    Let $p$ and $q$ be distinct odd \CrefAndHyperref{definition:prime_number}{prime numbers}. Then
    $$\left(\frac{p}{q}\right) \left(\frac{q}{p}\right) = (-1)^{\frac{p-1}{2} \cdot \frac{q-1}{2}}.$$
    Equivalently,
    $$\left(\frac{p}{q}\right) = \begin{cases}
        \left(\frac{q}{p}\right) & \text{if } p \equiv 1 \pmod{4} \text{ or } q \equiv 1 \pmod{4}, \\
        -\left(\frac{q}{p}\right) & \text{if } p \equiv q \equiv 3 \pmod{4}.
    \end{cases}$$
\end{theorem}

\begin{theorem} \label{theorem:wilsons_theorem}
    \TODO{AI generated, verify}
    Let $p$ be a \CrefAndHyperref{definition:prime_number}{prime number}. Then
    $$(p-1)! \equiv -1 \pmod{p}.$$
    Conversely, if $n \in \bbZ_{>1}$ satisfies $(n-1)! \equiv -1 \pmod{n}$, then $n$ is prime.
\end{theorem}



% ============================================================
% Section 10: Arithmetic Functions
% ============================================================
\section{Arithmetic Functions}

\begin{definition}[Arithmetic function] \label{definition:arithmetic_function_on_positive_integers_with_values_in_a_ring}
Let $R$ be a ring. An \hldef{arithmetic function on $\mathbb{Z}_{>0}$ with values in $R$} is a function $f: \mathbb{Z}_{>0} \to R$.
\end{definition}

\begin{definition}[Multiplicative function] \label{definition:multiplicative_function_on_positive_integers_with_values_in_a_ring}
Let $R$ be a ring. A \CrefAndHyperref{definition:arithmetic_function_on_positive_integers_with_values_in_a_ring}{arithmetic function on $\mathbb{Z}_{>0}$ with values in $R$} $f: \mathbb{Z}_{>0} \to R$ is called \hldef{multiplicative} if $f(1) = 1$ and for all $m, n \in \mathbb{Z}_{>0}$ with $\gcd(m, n) = 1$,
$$f(mn) = f(m)f(n).$$
\end{definition}

\begin{definition}[Completely multiplicative function] \label{definition:completely_multiplicative_function_on_positive_integers_with_values_in_a_ring}
Let $R$ be a ring. A \CrefAndHyperref{definition:arithmetic_function_on_positive_integers_with_values_in_a_ring}{arithmetic function on $\mathbb{Z}_{>0}$ with values in $R$} $f: \mathbb{Z}_{>0} \to R$ is called \hldef{completely multiplicative} if $f(1) = 1$ and for all $m, n \in \mathbb{Z}_{>0}$,
$$f(mn) = f(m)f(n).$$
\end{definition}

\begin{definition} \label{definition:dirichlet_convolution_of_arithmetic_functions}
    \TODO{AI generated, verify}
    Let $f, g: \bbZ_{>0} \to \bbC$ be \CrefAndHyperrefIfExist{definition:arithmetic_function_on_positive_integers_with_values_in_a_ring}{arithmetic functions}. Their \hldef{Dirichlet convolution}, denoted by \hl{$(f * g)$}, is the arithmetic function defined by
    $$\hlin{(f * g)(n) = \sum_{d \mid n} f(d)\, g\left(\frac{n}{d}\right)}$$
    for each $n \in \bbZ_{>0}$. Dirichlet convolution is commutative and associative, with identity element $\varepsilon$, where $\varepsilon(1) = 1$ and $\varepsilon(n) = 0$ for $n > 1$.
\end{definition}

\begin{definition}[Möbius function] \label{definition:mobius_function_on_positive_integers}
The \hldef{Möbius function}, denoted by \hl{$\mu$}, is the \CrefAndHyperref{definition:arithmetic_function_on_positive_integers_with_values_in_a_ring}{arithmetic function on $\mathbb{Z}_{>0}$ with values in $\mathbb{Z}$} defined by
$$\mu(1) = 1,$$
and for $n \in \mathbb{Z}_{>0}$ with $n > 1$,
$$\mu(n) = \begin{cases} (-1)^k & \text{if } n = p_1 p_2 \cdots p_k \text{ for distinct primes } p_1, \dots, p_k, \\ 0 & \text{otherwise.} \end{cases}$$
The Möbius function is a \CrefAndHyperref{definition:multiplicative_function_on_positive_integers_with_values_in_a_ring}{multiplicative function}.
\end{definition}

\begin{proposition} \label{proposition:mobius_inversion_formula}
    Let $R$ be a \CrefAndHyperrefIfExist{definition:ring}{ring}, and let $f, g: \bbZ_{>0} \to R$ be \CrefAndHyperrefIfExist{definition:arithmetic_function_on_positive_integers_with_values_in_a_ring}{arithmetic functions}. The following are equivalent:
    \begin{enumerate}
        \item For all $n \in \bbZ_{>0}$, $g(n) = \sum_{d \mid n} f(d).$
        \item For all $n \in \bbZ_{>0}$, $f(n) = \sum_{d \mid n} \mu\left(\frac{n}{d}\right) g(d),$ where $\mu$ is the \CrefAndHyperrefIfExist{definition:mobius_function_on_positive_integers}{Möbius function}, i.e. $f = \mu * g$\CrefIfExists{definition:dirichlet_convolution_of_arithmetic_functions}.
    \end{enumerate}
    Equivalently, $f(n) = \sum_{d \mid n} \mu(d)\, g\left(\frac{n}{d}\right).$
\end{proposition}



% ============================================================
% Section 11: Cyclotomic Polynomials
% ============================================================
\section{Cyclotomic Polynomials}

\begin{definition} \label{definition:cyclotomic_polynomial}
    Given an integer $n \geq 1$, the \hldef{$n$th cyclotomic polynomial} \hl{$\Phi_n(X)$} is the element of $\bbZ[X]$\CrefIfExists{definition:polynomial_ring_over_a_commutative_ring} defined inductively as follows:
    \begin{enumerate}
        \item $\Phi_1(X) = X-1$
        \item 
        $$\Phi_n(X) = \frac{X^n-1}{\prod_{\substack{d \mid n \\ d < n}} \Phi_d(X)}.$$
    \end{enumerate}

\end{definition}